% concatenated from neurips_2025.tex, math_commands.tex, tex/main.tex, tex/appendix.tex
\documentclass{article}


% if you need to pass options to natbib, use, e.g.:
%     \PassOptionsToPackage{numbers, compress}{natbib}
% before loading neurips_2025


% ready for submission
% \usepackage{neurips_2025}
% to compile a preprint version, e.g., for submission to arXiv, add add the
% [preprint] option:
%     \usepackage[preprint]{neurips_2025}


% to compile a camera-ready version, add the [final] option, e.g.:
\usepackage[final]{neurips_2025}


% to avoid loading the natbib package, add option nonatbib:
%    \usepackage[nonatbib]{neurips_2025}

\usepackage[utf8]{inputenc} % allow utf-8 input
\usepackage[T1]{fontenc}    % use 8-bit T1 fonts
\usepackage{hyperref}       % hyperlinks
\usepackage{url}            % simple URL typesetting
\usepackage{booktabs}       % professional-quality tables
\usepackage{amsfonts}       % blackboard math symbols
\usepackage{nicefrac}       % compact symbols for 1/2, etc.
\usepackage{microtype}      % microtypography
\usepackage{xcolor}         % colors

 
\usepackage{hyperref}
\usepackage{url} 
\usepackage{amsfonts} 
\usepackage{amssymb,amsmath,amsthm} 
\usepackage{mathabx}
\usepackage{mathrsfs}
\usepackage{xcolor}  
\usepackage{enumerate}
\usepackage{graphicx}
\usepackage{float}  
\usepackage{booktabs} 
\usepackage[nameinlink]{cleveref}
\usepackage[ruled,linesnumbered]{algorithm2e}
\crefname{algocf}{alg.}{algs.}
\Crefname{algocf}{Algorithm}{Algorithms}
\Crefname{equation}{Eq.}{Eqs.} 
\usepackage{attrib} 
\usepackage{subcaption}
 
%%%%% NEW MATH DEFINITIONS %%%%%

\usepackage{amsmath,amsfonts,bm}

% Mark sections of captions for referring to divisions of figures
\newcommand{\figleft}{{\em (Left)}}
\newcommand{\figcenter}{{\em (Center)}}
\newcommand{\figright}{{\em (Right)}}
\newcommand{\figtop}{{\em (Top)}}
\newcommand{\figbottom}{{\em (Bottom)}}
\newcommand{\captiona}{{\em (a)}}
\newcommand{\captionb}{{\em (b)}}
\newcommand{\captionc}{{\em (c)}}
\newcommand{\captiond}{{\em (d)}}

% Highlight a newly defined term
\newcommand{\newterm}[1]{{\bf #1}}


% Figure reference, lower-case.
\def\figref#1{figure~\ref{#1}}
% Figure reference, capital. For start of sentence
\def\Figref#1{Figure~\ref{#1}}
\def\twofigref#1#2{figures \ref{#1} and \ref{#2}}
\def\quadfigref#1#2#3#4{figures \ref{#1}, \ref{#2}, \ref{#3} and \ref{#4}}
% Section reference, lower-case.
\def\secref#1{section~\ref{#1}}
% Section reference, capital.
\def\Secref#1{Section~\ref{#1}}
% Reference to two sections.
\def\twosecrefs#1#2{sections \ref{#1} and \ref{#2}}
% Reference to three sections.
\def\secrefs#1#2#3{sections \ref{#1}, \ref{#2} and \ref{#3}}
% Reference to an equation, lower-case.
\def\eqref#1{equation~\ref{#1}}
% Reference to an equation, upper case
\def\Eqref#1{Equation~\ref{#1}}
% A raw reference to an equation---avoid using if possible
\def\plaineqref#1{\ref{#1}}
% Reference to a chapter, lower-case.
\def\chapref#1{chapter~\ref{#1}}
% Reference to an equation, upper case.
\def\Chapref#1{Chapter~\ref{#1}}
% Reference to a range of chapters
\def\rangechapref#1#2{chapters\ref{#1}--\ref{#2}}
% Reference to an algorithm, lower-case.
\def\algref#1{algorithm~\ref{#1}}
% Reference to an algorithm, upper case.
\def\Algref#1{Algorithm~\ref{#1}}
\def\twoalgref#1#2{algorithms \ref{#1} and \ref{#2}}
\def\Twoalgref#1#2{Algorithms \ref{#1} and \ref{#2}}
% Reference to a part, lower case
\def\partref#1{part~\ref{#1}}
% Reference to a part, upper case
\def\Partref#1{Part~\ref{#1}}
\def\twopartref#1#2{parts \ref{#1} and \ref{#2}}

\def\ceil#1{\lceil #1 \rceil}
\def\floor#1{\lfloor #1 \rfloor}
\def\1{\bm{1}}
\newcommand{\train}{\mathcal{D}}
\newcommand{\valid}{\mathcal{D_{\mathrm{valid}}}}
\newcommand{\test}{\mathcal{D_{\mathrm{test}}}}

\def\eps{{\epsilon}}


% Random variables
\def\reta{{\textnormal{$\eta$}}}
\def\ra{{\textnormal{a}}}
\def\rb{{\textnormal{b}}}
\def\rc{{\textnormal{c}}}
\def\rd{{\textnormal{d}}}
\def\re{{\textnormal{e}}}
\def\rf{{\textnormal{f}}}
\def\rg{{\textnormal{g}}}
\def\rh{{\textnormal{h}}}
\def\ri{{\textnormal{i}}}
\def\rj{{\textnormal{j}}}
\def\rk{{\textnormal{k}}}
\def\rl{{\textnormal{l}}}
% rm is already a command, just don't name any random variables m
\def\rn{{\textnormal{n}}}
\def\ro{{\textnormal{o}}}
\def\rp{{\textnormal{p}}}
\def\rq{{\textnormal{q}}}
\def\rr{{\textnormal{r}}}
\def\rs{{\textnormal{s}}}
\def\rt{{\textnormal{t}}}
\def\ru{{\textnormal{u}}}
\def\rv{{\textnormal{v}}}
\def\rw{{\textnormal{w}}}
\def\rx{{\textnormal{x}}}
\def\ry{{\textnormal{y}}}
\def\rz{{\textnormal{z}}}

% Random vectors
\def\rvepsilon{{\mathbf{\epsilon}}}
\def\rvtheta{{\mathbf{\theta}}}
\def\rva{{\mathbf{a}}}
\def\rvb{{\mathbf{b}}}
\def\rvc{{\mathbf{c}}}
\def\rvd{{\mathbf{d}}}
\def\rve{{\mathbf{e}}}
\def\rvf{{\mathbf{f}}}
\def\rvg{{\mathbf{g}}}
\def\rvh{{\mathbf{h}}}
\def\rvu{{\mathbf{i}}}
\def\rvj{{\mathbf{j}}}
\def\rvk{{\mathbf{k}}}
\def\rvl{{\mathbf{l}}}
\def\rvm{{\mathbf{m}}}
\def\rvn{{\mathbf{n}}}
\def\rvo{{\mathbf{o}}}
\def\rvp{{\mathbf{p}}}
\def\rvq{{\mathbf{q}}}
\def\rvr{{\mathbf{r}}}
\def\rvs{{\mathbf{s}}}
\def\rvt{{\mathbf{t}}}
\def\rvu{{\mathbf{u}}}
\def\rvv{{\mathbf{v}}}
\def\rvw{{\mathbf{w}}}
\def\rvx{{\mathbf{x}}}
\def\rvy{{\mathbf{y}}}
\def\rvz{{\mathbf{z}}}

% Elements of random vectors
\def\erva{{\textnormal{a}}}
\def\ervb{{\textnormal{b}}}
\def\ervc{{\textnormal{c}}}
\def\ervd{{\textnormal{d}}}
\def\erve{{\textnormal{e}}}
\def\ervf{{\textnormal{f}}}
\def\ervg{{\textnormal{g}}}
\def\ervh{{\textnormal{h}}}
\def\ervi{{\textnormal{i}}}
\def\ervj{{\textnormal{j}}}
\def\ervk{{\textnormal{k}}}
\def\ervl{{\textnormal{l}}}
\def\ervm{{\textnormal{m}}}
\def\ervn{{\textnormal{n}}}
\def\ervo{{\textnormal{o}}}
\def\ervp{{\textnormal{p}}}
\def\ervq{{\textnormal{q}}}
\def\ervr{{\textnormal{r}}}
\def\ervs{{\textnormal{s}}}
\def\ervt{{\textnormal{t}}}
\def\ervu{{\textnormal{u}}}
\def\ervv{{\textnormal{v}}}
\def\ervw{{\textnormal{w}}}
\def\ervx{{\textnormal{x}}}
\def\ervy{{\textnormal{y}}}
\def\ervz{{\textnormal{z}}}

% Random matrices
\def\rmA{{\mathbf{A}}}
\def\rmB{{\mathbf{B}}}
\def\rmC{{\mathbf{C}}}
\def\rmD{{\mathbf{D}}}
\def\rmE{{\mathbf{E}}}
\def\rmF{{\mathbf{F}}}
\def\rmG{{\mathbf{G}}}
\def\rmH{{\mathbf{H}}}
\def\rmI{{\mathbf{I}}}
\def\rmJ{{\mathbf{J}}}
\def\rmK{{\mathbf{K}}}
\def\rmL{{\mathbf{L}}}
\def\rmM{{\mathbf{M}}}
\def\rmN{{\mathbf{N}}}
\def\rmO{{\mathbf{O}}}
\def\rmP{{\mathbf{P}}}
\def\rmQ{{\mathbf{Q}}}
\def\rmR{{\mathbf{R}}}
\def\rmS{{\mathbf{S}}}
\def\rmT{{\mathbf{T}}}
\def\rmU{{\mathbf{U}}}
\def\rmV{{\mathbf{V}}}
\def\rmW{{\mathbf{W}}}
\def\rmX{{\mathbf{X}}}
\def\rmY{{\mathbf{Y}}}
\def\rmZ{{\mathbf{Z}}}

% Elements of random matrices
\def\ermA{{\textnormal{A}}}
\def\ermB{{\textnormal{B}}}
\def\ermC{{\textnormal{C}}}
\def\ermD{{\textnormal{D}}}
\def\ermE{{\textnormal{E}}}
\def\ermF{{\textnormal{F}}}
\def\ermG{{\textnormal{G}}}
\def\ermH{{\textnormal{H}}}
\def\ermI{{\textnormal{I}}}
\def\ermJ{{\textnormal{J}}}
\def\ermK{{\textnormal{K}}}
\def\ermL{{\textnormal{L}}}
\def\ermM{{\textnormal{M}}}
\def\ermN{{\textnormal{N}}}
\def\ermO{{\textnormal{O}}}
\def\ermP{{\textnormal{P}}}
\def\ermQ{{\textnormal{Q}}}
\def\ermR{{\textnormal{R}}}
\def\ermS{{\textnormal{S}}}
\def\ermT{{\textnormal{T}}}
\def\ermU{{\textnormal{U}}}
\def\ermV{{\textnormal{V}}}
\def\ermW{{\textnormal{W}}}
\def\ermX{{\textnormal{X}}}
\def\ermY{{\textnormal{Y}}}
\def\ermZ{{\textnormal{Z}}}

% Vectors
\def\vzero{{\bm{0}}}
\def\vone{{\bm{1}}}
\def\vmu{{\bm{\mu}}}
\def\vtheta{{\bm{\theta}}}
\def\va{{\bm{a}}}
\def\vb{{\bm{b}}}
\def\vc{{\bm{c}}}
\def\vd{{\bm{d}}}
\def\ve{{\bm{e}}}
\def\vf{{\bm{f}}}
\def\vg{{\bm{g}}}
\def\vh{{\bm{h}}}
\def\vi{{\bm{i}}}
\def\vj{{\bm{j}}}
\def\vk{{\bm{k}}}
\def\vl{{\bm{l}}}
\def\vm{{\bm{m}}}
\def\vn{{\bm{n}}}
\def\vo{{\bm{o}}}
\def\vp{{\bm{p}}}
\def\vq{{\bm{q}}}
\def\vr{{\bm{r}}}
\def\vs{{\bm{s}}}
\def\vt{{\bm{t}}}
\def\vu{{\bm{u}}}
\def\vv{{\bm{v}}}
\def\vw{{\bm{w}}}
\def\vx{{\bm{x}}}
\def\vy{{\bm{y}}}
\def\vz{{\bm{z}}}

% Elements of vectors
\def\evalpha{{\alpha}}
\def\evbeta{{\beta}}
\def\evepsilon{{\epsilon}}
\def\evlambda{{\lambda}}
\def\evomega{{\omega}}
\def\evmu{{\mu}}
\def\evpsi{{\psi}}
\def\evsigma{{\sigma}}
\def\evtheta{{\theta}}
\def\eva{{a}}
\def\evb{{b}}
\def\evc{{c}}
\def\evd{{d}}
\def\eve{{e}}
\def\evf{{f}}
\def\evg{{g}}
\def\evh{{h}}
\def\evi{{i}}
\def\evj{{j}}
\def\evk{{k}}
\def\evl{{l}}
\def\evm{{m}}
\def\evn{{n}}
\def\evo{{o}}
\def\evp{{p}}
\def\evq{{q}}
\def\evr{{r}}
\def\evs{{s}}
\def\evt{{t}}
\def\evu{{u}}
\def\evv{{v}}
\def\evw{{w}}
\def\evx{{x}}
\def\evy{{y}}
\def\evz{{z}}

% Matrix
\def\mA{{\bm{A}}}
\def\mB{{\bm{B}}}
\def\mC{{\bm{C}}}
\def\mD{{\bm{D}}}
\def\mE{{\bm{E}}}
\def\mF{{\bm{F}}}
\def\mG{{\bm{G}}}
\def\mH{{\bm{H}}}
\def\mI{{\bm{I}}}
\def\mJ{{\bm{J}}}
\def\mK{{\bm{K}}}
\def\mL{{\bm{L}}}
\def\mM{{\bm{M}}}
\def\mN{{\bm{N}}}
\def\mO{{\bm{O}}}
\def\mP{{\bm{P}}}
\def\mQ{{\bm{Q}}}
\def\mR{{\bm{R}}}
\def\mS{{\bm{S}}}
\def\mT{{\bm{T}}}
\def\mU{{\bm{U}}}
\def\mV{{\bm{V}}}
\def\mW{{\bm{W}}}
\def\mX{{\bm{X}}}
\def\mY{{\bm{Y}}}
\def\mZ{{\bm{Z}}}
\def\mBeta{{\bm{\beta}}}
\def\mPhi{{\bm{\Phi}}}
\def\mLambda{{\bm{\Lambda}}}
\def\mSigma{{\bm{\Sigma}}}

% Tensor
\DeclareMathAlphabet{\mathsfit}{\encodingdefault}{\sfdefault}{m}{sl}
\SetMathAlphabet{\mathsfit}{bold}{\encodingdefault}{\sfdefault}{bx}{n}
\newcommand{\tens}[1]{\bm{\mathsfit{#1}}}
\def\tA{{\tens{A}}}
\def\tB{{\tens{B}}}
\def\tC{{\tens{C}}}
\def\tD{{\tens{D}}}
\def\tE{{\tens{E}}}
\def\tF{{\tens{F}}}
\def\tG{{\tens{G}}}
\def\tH{{\tens{H}}}
\def\tI{{\tens{I}}}
\def\tJ{{\tens{J}}}
\def\tK{{\tens{K}}}
\def\tL{{\tens{L}}}
\def\tM{{\tens{M}}}
\def\tN{{\tens{N}}}
\def\tO{{\tens{O}}}
\def\tP{{\tens{P}}}
\def\tQ{{\tens{Q}}}
\def\tR{{\tens{R}}}
\def\tS{{\tens{S}}}
\def\tT{{\tens{T}}}
\def\tU{{\tens{U}}}
\def\tV{{\tens{V}}}
\def\tW{{\tens{W}}}
\def\tX{{\tens{X}}}
\def\tY{{\tens{Y}}}
\def\tZ{{\tens{Z}}}


% Graph
\def\gA{{\mathcal{A}}}
\def\gB{{\mathcal{B}}}
\def\gC{{\mathcal{C}}}
\def\gD{{\mathcal{D}}}
\def\gE{{\mathcal{E}}}
\def\gF{{\mathcal{F}}}
\def\gG{{\mathcal{G}}}
\def\gH{{\mathcal{H}}}
\def\gI{{\mathcal{I}}}
\def\gJ{{\mathcal{J}}}
\def\gK{{\mathcal{K}}}
\def\gL{{\mathcal{L}}}
\def\gM{{\mathcal{M}}}
\def\gN{{\mathcal{N}}}
\def\gO{{\mathcal{O}}}
\def\gP{{\mathcal{P}}}
\def\gQ{{\mathcal{Q}}}
\def\gR{{\mathcal{R}}}
\def\gS{{\mathcal{S}}}
\def\gT{{\mathcal{T}}}
\def\gU{{\mathcal{U}}}
\def\gV{{\mathcal{V}}}
\def\gW{{\mathcal{W}}}
\def\gX{{\mathcal{X}}}
\def\gY{{\mathcal{Y}}}
\def\gZ{{\mathcal{Z}}}

% Sets
\def\sA{{\mathbb{A}}}
\def\sB{{\mathbb{B}}}
\def\sC{{\mathbb{C}}}
\def\sD{{\mathbb{D}}}
% Don't use a set called E, because this would be the same as our symbol
% for expectation.
\def\sF{{\mathbb{F}}}
\def\sG{{\mathbb{G}}}
\def\sH{{\mathbb{H}}}
\def\sI{{\mathbb{I}}}
\def\sJ{{\mathbb{J}}}
\def\sK{{\mathbb{K}}}
\def\sL{{\mathbb{L}}}
\def\sM{{\mathbb{M}}}
\def\sN{{\mathbb{N}}}
\def\sO{{\mathbb{O}}}
\def\sP{{\mathbb{P}}}
\def\sQ{{\mathbb{Q}}}
\def\sR{{\mathbb{R}}}
\def\sS{{\mathbb{S}}}
\def\sT{{\mathbb{T}}}
\def\sU{{\mathbb{U}}}
\def\sV{{\mathbb{V}}}
\def\sW{{\mathbb{W}}}
\def\sX{{\mathbb{X}}}
\def\sY{{\mathbb{Y}}}
\def\sZ{{\mathbb{Z}}}

% Entries of a matrix
\def\emLambda{{\Lambda}}
\def\emA{{A}}
\def\emB{{B}}
\def\emC{{C}}
\def\emD{{D}}
\def\emE{{E}}
\def\emF{{F}}
\def\emG{{G}}
\def\emH{{H}}
\def\emI{{I}}
\def\emJ{{J}}
\def\emK{{K}}
\def\emL{{L}}
\def\emM{{M}}
\def\emN{{N}}
\def\emO{{O}}
\def\emP{{P}}
\def\emQ{{Q}}
\def\emR{{R}}
\def\emS{{S}}
\def\emT{{T}}
\def\emU{{U}}
\def\emV{{V}}
\def\emW{{W}}
\def\emX{{X}}
\def\emY{{Y}}
\def\emZ{{Z}}
\def\emSigma{{\Sigma}}

% entries of a tensor
% Same font as tensor, without \bm wrapper
\newcommand{\etens}[1]{\mathsfit{#1}}
\def\etLambda{{\etens{\Lambda}}}
\def\etA{{\etens{A}}}
\def\etB{{\etens{B}}}
\def\etC{{\etens{C}}}
\def\etD{{\etens{D}}}
\def\etE{{\etens{E}}}
\def\etF{{\etens{F}}}
\def\etG{{\etens{G}}}
\def\etH{{\etens{H}}}
\def\etI{{\etens{I}}}
\def\etJ{{\etens{J}}}
\def\etK{{\etens{K}}}
\def\etL{{\etens{L}}}
\def\etM{{\etens{M}}}
\def\etN{{\etens{N}}}
\def\etO{{\etens{O}}}
\def\etP{{\etens{P}}}
\def\etQ{{\etens{Q}}}
\def\etR{{\etens{R}}}
\def\etS{{\etens{S}}}
\def\etT{{\etens{T}}}
\def\etU{{\etens{U}}}
\def\etV{{\etens{V}}}
\def\etW{{\etens{W}}}
\def\etX{{\etens{X}}}
\def\etY{{\etens{Y}}}
\def\etZ{{\etens{Z}}}

% The true underlying data generating distribution
\newcommand{\pdata}{p_{\rm{data}}}
% The empirical distribution defined by the training set
\newcommand{\ptrain}{\hat{p}_{\rm{data}}}
\newcommand{\Ptrain}{\hat{P}_{\rm{data}}}
% The model distribution
\newcommand{\pmodel}{p_{\rm{model}}}
\newcommand{\Pmodel}{P_{\rm{model}}}
\newcommand{\ptildemodel}{\tilde{p}_{\rm{model}}}
% Stochastic autoencoder distributions
\newcommand{\pencode}{p_{\rm{encoder}}}
\newcommand{\pdecode}{p_{\rm{decoder}}}
\newcommand{\precons}{p_{\rm{reconstruct}}}

\newcommand{\laplace}{\mathrm{Laplace}} % Laplace distribution

\newcommand{\E}{\mathbb{E}}
\newcommand{\Ls}{\mathcal{L}}
\newcommand{\R}{\mathbb{R}}
\newcommand{\emp}{\tilde{p}}
\newcommand{\lr}{\alpha}
\newcommand{\reg}{\lambda}
\newcommand{\rect}{\mathrm{rectifier}}
\newcommand{\softmax}{\mathrm{softmax}}
\newcommand{\sigmoid}{\sigma}
\newcommand{\softplus}{\zeta}
\newcommand{\KL}{D_{\mathrm{KL}}}
\newcommand{\Var}{\mathrm{Var}}
\newcommand{\standarderror}{\mathrm{SE}}
\newcommand{\Cov}{\mathrm{Cov}}
% Wolfram Mathworld says $L^2$ is for function spaces and $\ell^2$ is for vectors
% But then they seem to use $L^2$ for vectors throughout the site, and so does
% wikipedia.
\newcommand{\normlzero}{L^0}
\newcommand{\normlone}{L^1}
\newcommand{\normltwo}{L^2}
\newcommand{\normlp}{L^p}
\newcommand{\normmax}{L^\infty}

\newcommand{\parents}{Pa} % See usage in notation.tex. Chosen to match Daphne's book.

\DeclareMathOperator*{\argmax}{arg\,max}
\DeclareMathOperator*{\argmin}{arg\,min}

\DeclareMathOperator{\sign}{sign}
\DeclareMathOperator{\Tr}{Tr}
\let\ab\allowbreak




\def\mathunderline#1#2{\color{#1}\underline{{\color{black}#2}}\color{black}}
\newcommand{\hlc}[2][yellow]{{%
		\colorlet{foo}{#1}%
		\sethlcolor{foo}\hl{#2}}%
}
\newcommand{\EE}{\mathbb{E}}
\newcommand{\FF}{\mathcal{F}}
\newcommand{\MM}{\mathscr{M}}
\newcommand{\calA}{\mathcal{A}}
\newcommand{\calF}{\mathcal{F}}
\newcommand{\calO}{\mathcal{O}}
\newcommand{\calP}{\mathcal{P}}
\newcommand{\calR}{\mathcal{R}}
\newcommand{\calS}{\mathcal{S}}
\newcommand{\calT}{\mathcal{T}}


\newcommand{\bias}{\mathfrak{b}} 
\newcommand{\MSE}{\mathsf{L}}
%\newcommand{\Var}{\mathfrak{V}}
\newcommand{\Leb}{\mathrm{Leb}}
%\newcommand{\Cov}{\mathrm{Cov}} 
\newcommand{\Trace}{\mathrm{Tr}} 
\newcommand{\Project}{\mathrm{Proj}} 
%%%%%%%%%%%%%%%%%%%%%%%%%%%%%%%%%%%%%%%%%%%%%%%%%%%%%%
 \newcommand{\bfQ}{\mathbf{Q}}
 \newcommand{\bfI}{\mathbf{I}}
 \newcommand{\bfP}{\mathbf{P}}
 \newcommand{\bfV}{\mathbf{V}}
 \newcommand{\bfD}{\mathbf{D}}
 \newcommand{\bfLambda}{\mathbf{\Lambda}}
 \newcommand{\bfDelta}{\mathbf{\Delta}}
 \newcommand{\fraki}{\mathfrak{i}}
 \newcommand{\frakj}{\mathfrak{j}}
 %%%%%%%%%%%%%%%%%%%%%%%%%%%%%%%%%%%%%%%%%%%%%%%%%%%%%%


\newcommand{\frakX}{\mathfrak{X}}


\newcommand{\PP}{\mathbb{P}}
\newcommand{\dd}{\mathrm{d}}
\newcommand{\RR}{\mathbb{R}}
\newcommand{\NN}{\mathbb{N}}
\newcommand{\TT}{\mathscr{T}}
\newcommand{\QQ}{\mathbf{Q}} 


\newcommand{\Idx}{\mathcal{I}}
\newcommand{\isom}{\mathrm{Isom}}
\newcommand{\aff}{\mathrm{Aff}} 
\newcommand{\Span}{\mathrm{Span}} 
\newcommand{\Proj}{\mathrm{Proj}} 
\newcommand{\prox}{\mathrm{prox}} 
\newcommand{\Image}{\mathrm{Im}} 
\newcommand{\Unif}{\operatorname{Uniform}} 

\makeatletter
\newcommand*\dotp{\mathpalette\dotp@{.5}}
\newcommand*\dotp@[2]{\mathbin{\vcenter{\hbox{\scalebox{#2}{$\m@th#1\bullet$}}}}}
\makeatother

\newcommand{\Lc}{\mathcal{L}^2_C}
\newcommand{\IndOne}{\boldsymbol{1}}
\newcommand\DistEq{\stackrel{D}{=}}
\newcommand\SimIID{\stackrel{iid}{\sim}}
\newcommand\ConvInProb{\xrightarrow{\ \PP \ }}
\newcommand\ConvAS{\xrightarrow{a.s.}}
%\newcommand\ConvWeak{\xrightarrow{\ w \ }}
\newcommand\ConvWeak{\xrightarrow{\ D \ }}
\newcommand\ConvInLp{\xrightarrow{L^p}}
\newcommand\Normal{N}

 
% BASIC MATH %

\newcommand{\bst}{ \hspace{1.5pt} | \hspace{1.5pt} }
\newcommand{\cst}{ \hspace{0.5pt} : \hspace{0.5pt} }
\newcommand{\sst}{ \hspace{2pt} ; \hspace{0.5pt} }
%\newcommand{\evl}{ \left| \right. }

\newcommand{\abs}[1]{\left| {#1} \right|}
%\newcommand{\ceil}[1]{\left\lceil #1 \right\rceil}
%\newcommand{\floor}[1]{\left\lfloor #1 \right\rfloor}

\newcommand{\dist}{\mathrm{dist}}

%\DeclareMathOperator*{\argmax}{arg\,max}
%\DeclareMathOperator*{\argmin}{arg\,min}

% Algebra
\newcommand{\permutation}{\mathrm{S}}
% Geometry
\newcommand{\sphere}{\mathbb{S}}


\renewcommand\labelitemii{$\circ$} 
\newcommand{\TODO}{{\color{red}TODO}}
\newcommand{\blue}[1]{{\color{blue} {#1}}}
\newcommand{\red}[1]{{\color{black} {#1}}}
\newcommand{\scrT}{\mathscr{T}}  
\newcommand{\const}[1]{{\color{blue}\mathit{C_{#1}}}}  
\usepackage[toc,page,header]{appendix}
\usepackage{minitoc}
\renewcommand \thepart{}
\renewcommand \partname{} 
\hypersetup{
	colorlinks,
	linkcolor={blue!75!black},
	citecolor={blue!75!black},
	urlcolor={blue!75!black}
}
 
\theoremstyle{definition}
\newtheorem{definition}{Definition}[section]
\newtheorem*{definition*}{Definition} 
\newtheorem{example}[definition]{Example}  
\theoremstyle{plain}
\newtheorem{theorem}[definition]{Theorem}
\newtheorem{conjecture}[definition]{Conjecture}
\newtheorem{lemma}[definition]{Lemma}
\newtheorem{proposition}[definition]{Proposition}
\newtheorem{corollary}[definition]{Corollary}
\newtheorem{assumption}[definition]{Assumption}
\newtheorem{condition}[definition]{Condition}
\newtheorem{exercise}[definition]{Exercise} 
\newtheorem*{theorem*}{Theorem}
\newtheorem*{lemma*}{Lemma}
\newtheorem*{proposition*}{Proposition}
\newtheorem{corollary*}{Corollary}   
\newtheorem{assumption*}{Assumption}
\newtheorem{condition*}{Condition}
\newtheorem{exercise*}{Exercise}
\newtheorem*{example*}{Example} 
\theoremstyle{remark}
\newtheorem*{remark}{Remark}
\usepackage{todonotes}
\usepackage{multirow}
\usepackage{makecell}
\usepackage{enumitem}
\usepackage[most]{tcolorbox}
\newtcolorbox{bluebox}{
  colback=blue!3!white,
  colframe=blue!60!black,
  % title=Research Question
}

\newcommand{\centerbox}[1]{\begin{center}\begin{minipage}{0.86\linewidth}
\begin{bluebox} 
\textit{{#1}}
\end{bluebox} \end{minipage}\end{center}}
\PassOptionsToPackage{numbers}{natbib}
\newcommand{\sfP}{\mathsf{P}} 

\title{On the Optimal Construction of Unbiased Gradient Estimators for Zeroth-Order Optimization}


% The \author macro works with any number of authors. There are two commands
% used to separate the names and addresses of multiple authors: \And and \AND.
%
% Using \And between authors leaves it to LaTeX to determine where to break the
% lines. Using \AND forces a line break at that point. So, if LaTeX puts 3 of 4
% authors names on the first line, and the last on the second line, try using
% \AND instead of \And before the third author name.


\author{%
  Shaocong Ma \\
  Department of Computer Science\\
  University of Maryland\\
  College Park, MD 20742, USA \\
  \texttt{scma0908@umd.edu} \\
  % examples of more authors
  \And
  Heng Huang\thanks{This work was partially supported by NSF IIS 2347592, 2348169, DBI 2405416, CCF 2348306, CNS 2347617, RISE 2536663.}\\
  Department of Computer Science\\
  University of Maryland\\
  College Park, MD 20742, USA \\
  \texttt{heng@umd.edu} \\
  % \AND
  % Coauthor \\
  % Affiliation \\
  % Address \\
  % \texttt{email} \\
  % \And
  % Coauthor \\
  % Affiliation \\
  % Address \\
  % \texttt{email} \\
  % \And
  % Coauthor \\
  % Affiliation \\
  % Address \\
  % \texttt{email} \\
}


\begin{document}


\doparttoc % Tell to minitoc to generate a toc for the parts
\faketableofcontents % Run a fake tableofcontents command for the partocs  


\maketitle


\begin{abstract}
  Zeroth-order optimization (ZOO) is an important framework for stochastic optimization when gradients are unavailable or expensive to compute. A potential limitation of existing ZOO methods is the bias inherent in most gradient estimators unless the perturbation stepsize vanishes. In this paper, we overcome this biasedness issue by proposing a novel family of \textit{unbiased} gradient estimators based solely on function evaluations. By reformulating directional derivatives as a telescoping series and sampling from carefully designed distributions, we construct estimators that eliminate bias while maintaining favorable variance. We analyze their theoretical properties, derive optimal scaling distributions and perturbation stepsizes of four specific constructions, and prove that SGD using the proposed estimators achieves optimal complexity for smooth non-convex objectives. Experiments on synthetic tasks and language model fine-tuning confirm the superior accuracy and convergence of our approach compared to standard methods.
\end{abstract}

\section{Introduction}
In this paper, we consider the problem of \textit{zeroth-order optimization (ZOO)}, where our goal is to solve the following  stochastic optimization problem:
\begin{align}\label{eq:optimization-problem}
\min_{x \in \RR^d} f(x) := \EE_{\xi \sim \Xi} f(x;\xi),
\end{align}
where $f(x;\xi)$ is a smooth loss function evaluated on data  $\xi$ drawn from a distribution $\Xi$.  In many practical scenarios, gradient information is either unavailable or prohibitively expensive to compute. Due to its versatility, ZOO has been widely adopted across various domains, including black-box adversarial attacks on machine learning models \citep{chen2017zoo, kurakin2016adversarial, papernot2017practical, cai2021zeroth, zhao2020towards}, physics-informed neural networks interfacing with external PDE solvers \citep{shen2024memory, Ma2025}, and reinforcement learning \citep{choromanski2018structured, lei2022zeroth, suh2022differentiable}. Recent research on ZOO also focuses on enhancing memory efficiency \citep{cai2022one,cai2022zeroth,li2024addax,sugiura2025elasticzo}, motivated in large part by fine-tuning large language models \citep{malladi2023fine, zhang2024revisiting, gautam2024variance, tang2024private, wang2024simultaneous, wang2025zo2}. 

Unlike first-order methods that rely on stochastic gradients $\nabla f(x; \xi)$, ZOO uses only function evaluations, without access to gradient information. To approximate gradients, several estimators have been proposed, including the one-point estimate $\hat{\nabla} f(x;\xi)= \frac{f(x+\mu v;\xi)}{\mu} v$ \citep{flaxman2005online,shamir2013complexity,bach2016highly,nesterov2017random,berahas2022theoretical} and two-point estimator $\hat{\nabla} f(x;\xi)= \frac{f(x+\mu v;\xi) - f(x;\xi)}{\mu}v$ \citep{ghadimi2013stochastic,duchi2015optimal,nesterov2017random} (see \Cref{appendix:gradient-estimators} for further discussions). The random direction $v$ is typically drawn from a Gaussian or uniform spherical distribution, while alternative choices have also gained increasing attention in recent years \citep{ghadimi2013stochastic, duchi2015optimal, ji2019improved, sahu2019towards, coope2020gradient, kozak2023zeroth, rando2023optimal, rando2024stochastic, ma2025revisiting,mi2025kerzoo}.



However, despite these advancements, a critical limitation arises in zeroth-order gradient estimation; that is, all widely used gradient estimators exhibit inherent bias. Specifically, unless the perturbation step size $\mu$  asymptotically tends to zero, these estimators yield persistently biased approximations of the true gradient. This inherent bias motivates a central question explored in this paper:

\centerbox{\textbf{Q1:} Is it possible to design an unbiased zeroth-order gradient estimator using only function evaluations?}  
 
\textbf{Contribution 1}: In this paper, we answer \textit{\textbf{Q1}} affirmatively. Contrary to the belief that zeroth-order gradient estimators must inherently be biased due to finite-step perturbations, we demonstrate that it is indeed possible to construct \textit{unbiased} gradient estimators using only function evaluations. Our key idea is to express $\nabla_v f(x)$ (in the deterministic setting), the directional derivative along the direction $v$, as a telescoping series:
\begin{align}\label{eq:telescoping}
    &\nabla_vf(x):= \lim_{\mu_n \to 0} \frac{f(x+\mu_n v) - f(x)}{\mu_n} \nonumber\\
    &=\sum_{n=1}^\infty p_n\left[\frac{f(x+\mu_1 v) - f(x)}{\mu_1} + \frac{1}{p_n}\left( \frac{f(x+\mu_{n+1} v) - f(x)}{\mu_{n+1}} - \frac{f(x+\mu_n v) - f(x)}{\mu_n}\right) \right], \\
    &\overset{(i)}{=} \EE_{n\sim \{p_n\}_{n=1}^\infty} \left[\frac{f(x+\mu_1 v) - f(x)}{\mu_1} + \frac{1}{p_n}\left( \frac{f(x+\mu_{n+1} v) - f(x)}{\mu_{n+1}} - \frac{f(x+\mu_n v) - f(x)}{\mu_n}\right) \right]\nonumber
\end{align} 
where the perturbation stepsize $\mu_n \to 0$ as $n \to \infty$, the sampling distribution $\{p_n\}_{n=1}^\infty$ form a probability distribution (that is, $0<p_i<1$ for all $i\in \NN$ and $\sum_{i=1}^\infty p_i = 1$), and the expectation representation (i) holds under mild regularity conditions (\Cref{prop:abs-conv}).   This formulation allows us to reinterpret the directional derivative as an expectation over $n \sim \{p_n\}_{n=1}^\infty$, enabling the construction of a unbiased gradient estimator family $\mathscr{P}$ (\Cref{def:unbiased-family}): 
\begin{align*}
    \hat{\nabla}_v f(x):= \EE_{n\sim p} \sfP(n,v), 
\end{align*} 
where $\sfP(n,v)$ is an unbiased estimator of $$\frac{f(x+\mu_1 v) - f(x)}{\mu_1} + \frac{1}{p_n}\left( \frac{f(x+\mu_{n+1} v) - f(x)}{\mu_{n+1}} - \frac{f(x+\mu_n v) - f(x)}{\mu_n}\right).$$ 
Within this framework, we propose four specific estimators, denoted as \textit{$\sfP_k$-estimator} for $k = 1, 2, 3, 4$, corresponding to the number of function evaluations required in each estimation. To the best of our knowledge, unbiased zeroth-order gradient estimators have received little attention in prior literature.  The only existing work we are aware of is the four-point estimator proposed by \citet{chen2020unbiased}, which shares the same telescoping structure and can be viewed as a special case of our $\sfP_4$-estimator.  

\textbf{Contribution 2:}  
Building on our unbiased estimator construction, we conduct a rigorous variance analysis on our proposed $\sfP_k$-estimators.  We first present a negative result for the $\sfP_1$-estimator;  although it requires fewer function evaluations, it may exhibit infinite variance under certain conditions (\Cref{thm:variance} (a)), which aligns with the one-point estimator in the randomized smoothing \citep{flaxman2005online}.   Next, we characterize the relation among the variance of  the $\sfP_k$-estimator ($k=2,3,4$), the perturbation stepsize sequence $\{\mu_n\}_{n=1}^\infty$, and the sampling distribution $\{p_n\}_{n=1}^\infty$ (\Cref{thm:variance} (b)). Identifying the optimal  choice of  $\{\mu_n\}_{n=1}^\infty$ and $\{p_n\}_{n=1}^\infty$ leads us to the following \textit{non-convex functional optimization} problem: 
\begin{align}\label{eq:functional-opt}
    \min_{\{\mu_n\}_{n=1}^\infty, \{p_n\}_{n=1}^\infty} \quad &\EE_{n \sim \{p_n\}_{n=1}^\infty} \left( \frac{\mu_n - \mu_{n+1}}{p_n}\right)^2  \\
    \text{subject to} \quad &0 < p_n < 1; \; \sum_{n=1}^\infty p_n=1; \sum_{n=1}^\infty \mu_n<\infty. \nonumber
\end{align} 
We present an explicit analytical solution to this optimization problem (\Cref{prop:lower_bound_varrho}), which reveals two key insights: (1) our constructed unbiased gradient estimators can achieve the same variance as the classical two-point estimator without introducing additional bias, leading to the best-possible complexity for SGD algorithm (\Cref{cor:sgd-convergence}); (2) a broad class of sampling distributions can achieve the minimum variance, extending beyond the specific choices considered in prior work \citep{chen2020unbiased}. While our theoretical results establish strong guarantees, an important practical question remains:

\centerbox{\textbf{Q2:} Given the optimal choice of $\{\mu_n\}_{n=1}^\infty$ and $\{p_n\}_{n=1}^\infty$, do the proposed unbiased estimators empirically outperform existing zeroth-order methods? }  
\textbf{Contribution 3:}  
To address \textit{\textbf{Q2}}, we empirically validate our proposed approach across both synthetic and practical tasks.   On estimating the gradient of mean-square and logistic losses, our method achieves significantly lower gradient estimation error compared to standard zeroth-order methods (\Cref{sec:synthetic}). Furthermore, when applied to fine-tuning large language models, the proposed estimators demonstrate faster convergence and higher final accuracy under the same number of function evaluations (\Cref{sec:llm}). These results confirm the practical advantages of our unbiased construction and underscore its effectiveness in modern zeroth-order optimization tasks.
  
\section{The Derivation of Unbiased Zeroth-Order Estimators}\label{sec:min-var} 

We will start from the deterministic case then turn to the stochastic case in \Cref{sec:convergence}. In this section, we formally derive a class of unbiased estimators for approximating the gradient $\nabla f(x)$ using only function evaluations. We also provide a sufficient condition under which the telescoping series in \Cref{eq:series} admits the expectation representation. All proofs are provided in the appendix.
% Our telescoping series coincide with the unbiased four-point estimator~\citep{chen2020unbiased}, while we go further in this direction to derive the sufficient condition on when such telescoping series can form an expectation, which significantly generalizes the result presented by \citet{chen2020unbiased}, where only the geometric distribution is used.
 

\subsection{Telescoping Series and Expectation Representation} 
For a fixed direction \(v\in\mathbb{R}^d\), the directional derivative of a differentiable function \(f:\mathbb{R}^d\to\mathbb{R}\) at \(x\) along the direction $v$ is defined as $$\nabla_v f(x) = \lim_{\mu\to 0} \frac{f(x+\mu v)-f(x)}{\mu}.$$ 
Then for any decreasing sequence \(\{\mu_n\}_{n=1}^\infty\) with \(\lim_{n \to \infty} \mu_n = 0\), one can express this directional derivative as the limit of a convergent sequence $\left\{ \frac{f(x+\mu_n v) - f(x)}{\mu_n} \right\}$: $$\nabla_v f(x) = \lim_{n \to \infty} \frac{f(x+\mu_n v) - f(x)}{\mu_n} . $$

This convergent sequence canonically induces a telescoping series with the same limit:
\begin{align}\label{eq:series}
    \nabla_v f(x) =\frac{f(x+\mu_1 v)-f(x)}{\mu_1}  + \sum_{n=1}^\infty \left[ \frac{f(x+\mu_{n+1} v)-f(x)}{\mu_{n+1}} - \frac{f(x+\mu_n v)-f(x)}{\mu_n} \right].
\end{align}

Next, consider a probability mass function (PMF) \(\{p_n\}_{n=1}^\infty\) with \(p_n > 0\) for all \(n\) and \(\sum_{n=1}^\infty p_n = 1\). When the series in \Cref{eq:series} is \textbf{\textit{absolutely convergent}}\footnote{We follow the standard definition from \citet{spivak2008calculus}: A series $\sum_{n=1}^\infty a_n$ is called \textit{convergent}, if the limit of its finite sum $\lim_{N\to \infty }\sum_{n=1}^N a_n$ exists. A series $\sum_{n=1}^\infty a_n$ is called \textit{absolutely convergent}, if the series $\sum_{n=1}^\infty |a_n|$ is convergent. See the formal definition in \Cref{appendix:absolute-convergence}.}, we can interpret it as an expectation over a discrete random variable $n$. That is, 
\begin{align}
\nabla_v f(x) &= \sum_{n=1}^\infty p_n\left[\frac{f(x+\mu_1 v) - f(x)}{\mu_1} + \frac{1}{p_n}\left( \frac{f(x+\mu_{n+1} v) - f(x)}{\mu_{n+1}} - \frac{f(x+\mu_n v) - f(x)}{\mu_n}\right) \right] \nonumber\\
&= \mathbb{E}\left[\frac{f(x+\mu_1 v) - f(x)}{\mu_1} + \frac{1}{p_n}\left( \frac{f(x+\mu_{n+1} v) - f(x)}{\mu_{n+1}} - \frac{f(x+\mu_n v) - f(x)}{\mu_n}\right) \right]. \label{eq:expectation}
\end{align}
% In particular, if we construct an estimator \(g\) such that $\mathbb{E}[g] = \nabla_v f(x),$ 
% then the product \(g\,v\) yields an unbiased estimator for the full gradient, \(\nabla f(x)\), since \(\mathbb{E}[g\,v] = \mathbb{E}[g \mid v]\,v = \nabla_v f(x)\,v = \nabla f(x)\).

\paragraph{On the Role of Absolute Convergence.}
The absolute convergence of the series in \Cref{eq:series} plays a critical role in interpreting the telescoping series as an expectation. This is due to the difference between the series convergence and the existence of expectation:  
\begin{itemize}
    \item \textit{The series convergence}: Consider the convergent series $\sum_{i=1}^\infty p_i x_i$. To evaluate its value, we can calculate the finite-sum $S_n:=\sum_{i=1}^n p_i x_i$; then we have $\sum_{i=1}^\infty p_i x_i= \lim_{n\to \infty} S_n$.
    \item \textit{The existence of expectation}:  Consider the random variable $X$ with $\PP(X=x_i)=p_i$ for $i \in \NN$. Its expectation $\EE[X]$ is also written as $\sum_{i=1}^\infty p_i x_i$. However, the notion of expectation must be well-defined independently of any ordering of outcomes. That is, for an arbitrary permutation $\sigma: \NN \to \NN$, all series $\sum_{i=1}^\infty p_{\sigma(i)} x_{\sigma(i)}$ should represent the same value $\EE[X]$.
\end{itemize}
As a result, a convergent series can yield different values depending on the order of summation (this result is called the Riemann series theorem \citep{riemann1868darstellbarkeit,spivak2008calculus}); however, the outcomes of a random variable requires a random variable's expectation to be well-defined regardless of any such ordering. While the expectation representation has been discussed in prior work (e.g., \citep{chen2020unbiased}), the lack of attention to absolute convergence has left the conditions ensuring unbiasedness underexplored. 

Due to this reason, we provide the following (mild) sufficient condition for ensuring the absolute convergence with adding a slightly stronger requirement on the objective function $f:\RR^d \to \RR$ and the sequence $\{ \mu_n\}_{n=1}^\infty$:
\begin{proposition}\label{prop:abs-conv}
    If the second-order continuously differentiable function $f:\RR^d \to \RR$ has $L$-Lipschitz continuous gradient and $\sum_{n=1}^\infty \mu_n < \infty$, then the series $$\sum_{n=1}^\infty p_n\left[\frac{f(x+\mu_1 v) - f(x)}{\mu_1} + \frac{1}{p_n}\left( \frac{f(x+\mu_{n+1} v) - f(x)}{\mu_{n+1}} - \frac{f(x+\mu_n v) - f(x)}{\mu_n}\right) \right]$$ is absolutely convergent and its limit is $\nabla_v f(x)$.
\end{proposition}
% \begin{proof}
%     The proof simply follows by expanding $f(x+\mu v)$ using Taylor's theorem at $x$; each entry of the series is upper bounded by $p_nC + (\mu_{n+1} + \mu_n)C'$ for some positive constant $C$ and $C'$. The proof details are postponed to \Cref{prop:abs-conv-appendix}.
% \end{proof} 

\subsection{The Construction of Unbiased Estimators}
With the expectation representation in place, we are ready to define the class of unbiased estimators explicitly.  
\begin{definition} \label{def:unbiased-family}
Suppose that the function $f:\RR^d\to \RR$ is  continuously differentiable and  $\{\mu_n\}_{n\ge1}$ is a positive sequence with $\lim_{n\to\infty}\mu_n=0$ such that the telescoping series
\[
\frac{f(x+\mu_1 v)-f(x)}{\mu_1} + \sum_{n=1}^\infty \left[ \frac{f(x+\mu_{n+1} v)-f(x)}{\mu_{n+1}}-\frac{f(x+\mu_n v)-f(x)}{\mu_n}\right]
\]
is absolutely convergent, the sequence $\{p_n\}_{n=1}^\infty$ forms a PMF, and $V$ is the distribution over $\RR^d$. Then the family of estimators $\mathscr{P}:=\mathscr{P}(f, \{ \mu_n\}_{n=1}^\infty, \{ p_n\}_{n=1}^\infty,V)$ denote the class of random variables such that for every $\mathsf{P}(n,v) \in \mathscr{P}$, it satisfies
    $$\EE[\mathsf{P}(n,v)\mid n,v]= \frac{f(x+\mu_1 v) - f(x)}{\mu_1} + \frac{1}{p_n}\left( \frac{f(x+\mu_{n+1} v) - f(x)}{\mu_{n+1}} - \frac{f(x+\mu_n v) - f(x)}{\mu_n}\right),$$
    where $v$ is sampled from $V$,  independent with $n\sim \{ p_n\}_{n=1}^\infty$. 
\end{definition} 
% \begin{remark}
% This class of estimators follow the construction of four-point unbiased gradient estimator proposed by \citet{chen2020unbiased}. Here, we significantly generalize it to include estimators that require less function evaluations. 
% \end{remark}

In the following theorem, we formally prove that our proposed class $\mathscr{P}$ is exactly the unbiased estimator of the gradient $\nabla f(x)$. 
\begin{theorem}[Unbiasedness]\label{thm:unbiased}
Let $\mathscr{P}:=\mathscr{P}(f, \{ \mu_n\}_{n=1}^\infty, \{ p_n\}_{n=1}^\infty,V)$ is defined as \Cref{def:unbiased-family}. Then, for any estimator $\mathsf{P}(n,v) \in \mathscr{P}$, the following hold:
\begin{enumerate}[label=(\alph*)]
    \item \(\mathbb{E} [\mathsf{P}(n,v)\mid v] = \nabla_v f(x)\); that is, \(\mathsf{P}(n,v)\) is an unbiased estimator of the directional derivative $\nabla_v f(x)$.
    \item If the random direction \(v\) is chosen independently of the sampling $n\sim \{p_n\}_{n=1}^\infty $ and satisfies \(\mathbb{E}[v\,v^\top]=I\), then
    \[
    \mathbb{E}_{n\sim \{p_n\}_{n=1}^\infty , v\sim V}\Bigl[\mathsf{P}(n,v)\,v\Bigr] = \nabla f(x),
    \]
    so that \(\mathsf{P}(n,v)\,v\) is an unbiased estimator of the gradient $ \nabla f(x)$.
\end{enumerate}
\end{theorem}
% \begin{proof}
% The proof directly follows the expectation representation constructed in \Cref{eq:expectation} and the definition of the directional derivative. We defer the full details to  \Cref{sec:unbiasedness}.
% \end{proof}  

\subsection{Specific Constructions}
In this subsection, we propose four concrete  constructions from the estimator family (\Cref{def:unbiased-family}) $$\mathscr{P}:=\mathscr{P}(f, \{ \mu_n\}_{n=1}^\infty, \{ p_n\}_{n=1}^\infty,V)$$ based on the number of function evaluations used in estimating the gradient. These constructions are designed to explore two main aspects: (1) the trade-off between the estimator variance and the number of function evaluations, allowing flexibility depending on the computational budget; and (2) a fundamental question purely driven by the theoretical interest:  \textit{What is the minimum number of function evaluations required to construct an unbiased gradient estimator?}

\paragraph{\(\mathsf{P}_4\)-Estimator.}  
This estimator corresponds to the four-point estimator originally proposed by \citet{chen2020unbiased} with slightly generalizing the choice of the perturbation stepsize sequence $\{\mu_n\}_{n=1}^\infty$ and the sampling distribution $\{p_n\}_{n=1}^\infty$%\footnote{In \citet{chen2020unbiased}'s work, only the exponential decaying $\mu_n := \mu_0 a^n$ and the geometric scaling distribution $p_n \propto c^n$ are considered.}
. For a given direction $v\sim V$ and $n\sim \{p_n\}_{n=1}^\infty$, the \(\mathsf{P}_4\)-estimator is defined as 
\begin{equation}
\mathsf{P}_4(n, v)= \frac{f(x+\mu_1 v)-f(x)}{\mu_1} + \frac{1}{p_n}\Biggl[ \frac{f(x+\mu_{n+1} v)-f(x)}{\mu_{n+1}} - \frac{f(x+\mu_n v)-f(x)}{\mu_n} \Biggr]. \label{eq:P4_def}
\end{equation}
This construction requires four function evaluations at: $x$, $x+\mu_1 v$, $x+\mu_n v$, and $x+\mu_{n+1}v$, exhibiting the  lowest variance and the most function evaluation counts among all members of $\mathscr{P}$. 


\paragraph{\(\mathsf{P}_3\)-Estimator.}  
We can reduce one function evaluation by introducing a selection random variable $\mathrm{U}_2\sim  \operatorname{Uniform}\left( \{ 0, 1 \}\right)$\footnote{Here we use $\operatorname{Uniform}(A)$ to represent the uniform distribution over the finite or compact set $A$.}. The estimator is then defined as
\begin{equation}
\mathsf{P}_3(n, v)= \frac{f(x+\mu_1 v)-f(x)}{\mu_1}\mathrm{U}_2 + \frac{1}{p_n}\Biggl[ \frac{f(x+\mu_{n+1} v)-f(x)}{\mu_{n+1}} - \frac{f(x+\mu_n v)-f(x)}{\mu_n} \Biggr](1-\mathrm{U}_2). \label{eq:P3_def}
\end{equation}
This construction randomly selects one of two pathways:  With probability \(1/2\), it uses the first term only and requires two function evaluations at \(x+\mu_1 v\) and \(x\); otherwise, it uses the second term and requires three function evaluations at \(x+\mu_n v\), \(x+\mu_{n+1} v\), and \(x\).  This estimator maintains unbiasedness as \(\mathsf{P}_4\), with slightly higher variance. 

 
\paragraph{\(\mathsf{P}_1\)- \& \(\mathsf{P}_2\)-Estimator.} The selection random variable can be naturally extended to construct \(\mathsf{P}_1\)- and \(\mathsf{P}_2\)-estimators as follows:
\begin{align}
\mathsf{P}_2(n, v)= &\frac{f(x+\mu_1 v)-f(x)}{\mu_1}\mathbb{I}_{\{\mathrm{U}_3=0\}} \label{eq:P2_def}\\ 
&+ \frac{1}{p_n}\Biggl[ \frac{f(x+\mu_{n+1} v)-f(x)}{\mu_{n+1}}\mathbb{I}_{\{\mathrm{U}_3=1\}} - \frac{f(x+\mu_n v)-f(x)}{\mu_n}\mathbb{I}_{\{\mathrm{U}_3=2\}} \Biggr] , \nonumber \\
\mathsf{P}_1(n, v)= &\frac{f(x+\mu_1 v)\mathbb{I}_{\{\mathrm{U}_4=1\}}-f(x)\mathbb{I}_{\{\mathrm{U}_4=0\}}}{\mu_1} \label{eq:P1_def}\\ 
&+ \frac{1}{p_n}\Biggl[ \frac{f(x+\mu_{n+1} v)\mathbb{I}_{\{\mathrm{U}_4=2\}}-f(x)\mathbb{I}_{\{\mathrm{U}_4=0\}}}{\mu_{n+1}}  - \frac{f(x+\mu_n v)\mathbb{I}_{\{\mathrm{U}_4=3\}}-f(x)\mathbb{I}_{\{\mathrm{U}_4=0\}}}{\mu_n}  \Biggr] , \nonumber
\end{align} 
where $\mathrm{U}_3 \sim \operatorname{Uniform}\left( \{ 0, 1,2 \}\right)$, $\mathrm{U}_4 \sim \operatorname{Uniform}\left( \{ 0, 1,2,3 \}\right)$, and $\mathbb{I}_{A}$ is the indicator function, which equals $1$ if the event $A$ occurs, and $0$ otherwise. of the event $A$. Remarkably, the construction of $\sfP_1$-estimator achieves unbiasedness using only a single function evaluation. However, we will show that in the next section,$\sfP_1$-estimator will have infinite variance under certain condition.  

 
\section{Variance Analysis of Unbiased Zeroth-Order Estimators}
% \input{tex/variance-table} 

In this section, we provide a theoretical analysis of the variance behavior for the unbiased estimator family $\mathscr{P} = \mathscr{P}(f, \{ \mu_n\}_{n=1}^\infty, \{ p_n\}_{n=1}^\infty, V)$ (\Cref{def:unbiased-family}). While the unbiasedness has been shown in \Cref{thm:unbiased}, their variances can differ dramatically depending on the estimator construction. In particular, we prove that the variance becomes unbounded (i.e., infinite) for certain constructions such as \(\sfP_1\)-estimator. We also provide finite-variance bounds for \(\sfP_k\)-estimators (for $k=2,3,4$) with matching the optimal variance under specific choices of $\{p_n\}$ and $\{\mu_n\}$. 
  
\subsection{Theoretical Analysis}
In the following result, we adopt the same condition as \Cref{prop:abs-conv} to ensure the expectation representation. 
\begin{theorem}\label{thm:variance}
Let $\mathscr{P}:=\mathscr{P}(f, \{ \mu_n\}_{n=1}^\infty, \{ p_n\}_{n=1}^\infty,V)$ is defined as \Cref{def:unbiased-family}. Suppose that $f:\RR^d \to \RR$ is second-order continuously differentiable and has $L$-Lipschitz continuous gradient, $\sum_{n=1}^\infty \mu_n < \infty$, and $V$ is the uniform distribution over the  sphere with the radius $\sqrt{d}$. Define 
\begin{align*}
    \mu:=\mu_1, \qquad\varrho:= \sum_{n=1}^\infty \frac{(\mu_{n+1} - \mu_{n})^2}{p_n}, \qquad \text{and} \qquad \varphi:=\sum_{n=1}^\infty  \frac{\mu_{n}^2}{p_n}.
\end{align*}
Then the following statements hold:
    \begin{enumerate}[label=(\alph*)]
        \item If there exists a point $x\in \RR^d$ such that the Hessian $\nabla^2 f(x)$ is positive definite and $f(x)\neq 0$, then the variances of the $\sfP_1$ for estimating $\nabla f(x)$ is infinite.  
        \item The variance of $\sfP_k$-estimator $\sfP_k(n,v) v$ ($k=2,3,4$) for estimating $\nabla f(x)$ is given by
        \begin{align*}
        \Var[\sfP_2(n,v) \, v ]  
    &\leq  \Var[\sfP_4(n,v)\, v] + \frac{L^2 }{3} d^3\mu^2 + \frac{L^2}{12}  d^3  \varrho  + \frac{L^2}{3} d^3   \varphi.\\
        \Var[\sfP_3(n,v) v] &\leq \Var[\sfP_4(n,v)\, v] + \frac{L^2 }{8} d^3\mu^2 + \frac{L^2}{8}  d^3 \varrho. \\ 
            \Var[\sfP_4(n,v) v] &\leq  (d-1) \| \nabla f(x) \|^2 +  \frac{3L^2}{4} d^3 \mu^2  +  \frac{L^2 d^3}{2}  \varrho.
        \end{align*} 
    \end{enumerate}
\end{theorem}
\begin{proof}
    Part (a) directly follows by analyzing the tail of $\frac{1}{p_n}\frac{f(x+\mu_n v)}{\mu_n}$  and leveraging the curvature from a positive definite Hessian. For the part (b), we simply decompose the variance of $\sfP_2(n,v) v$ and $\sfP_3(n,v) v$ into the variance of estimating $\sfP_4(n,v) v$ using 
\begin{align*}
    \Var[\sfP \, v ]  =     d\EE[(\sfP- \sfP_4(n,v))^2  ] + \Var[\sfP_4(n,v)\, v]
\end{align*}
for arbitrary $\sfP:=\sfP(n,v)\in \mathscr{P}$. 
    Then we apply the second-order Taylor expansions with the mean value theorem to control the finite-difference noise.  Full details and auxiliary lemmas are provided in \Cref{appendix:variance}. 
\end{proof}  
 
\paragraph{Comparison with Existing Literature.}  
\Cref{thm:variance} reveals that while $\mathsf{P}_1$ is unbiased, its variance can be infinite under certain conditions, making them unsuitable for SGD. In contrast, $\mathsf{P}_k$-estimator ($k=2,3,4$) offer the finite variance when $\{ \mu_n\}_{n=1}^\infty$ and $\{ p_n\}_{n=1}^\infty$ are appropriately selected. We will show it later that under the optimal setting, their variances match the optimal order of classical two-point estimators \citep{nesterov2017random} but with zero bias:  
$$\Var[\sfP_k \ v] = \calO(d\|\nabla f(x)\|^2 + d^3 \mu^2).$$
This variance will lead to the optimal function query complexity $\calO(\frac{d}{\epsilon^4})$ for achieving  $\epsilon$-accuracy in the gradient norm $\| \nabla f(x)\|$ \citep{duchi2015optimal}.


\paragraph{Comparison with the Noisy Oracle Setup}  In our work, we consider the exact function evaluation setting with noiseless values. In this case, our variance scales as $d^3 \mu^2$, which is worse than the $d^2 \mu^2$ of some specific biased estimators, which is mitigated by choosing a small enough $\mu$; the overall sample complexity remains optimal. However, in the noisy function evaluation setting, where each function evaluation may return a noisy value, a smaller $\mu$ amplifies the noise, leading to degraded performance. Several recent works have provided more refined analysis under noisy setups with improved variance behavior. Notably,  \citet{akhavan2024gradient} demonstrated that for highly smooth functions, the $\ell_1$ -randomization can reduce the variance scaling to $d^2 \mu^2$ with achieving the improved performance for highly smooth objective functions, which extends the existing $\ell_1$-randomization proposed by \citet{akhavan2022gradient}. Earlier work by \citet{gasnikov2017stochastic} analyzed the variance behavior in single-point and multi-point bandit feedback settings, and more recent developments further explore the impact of first-order smoothness in noisy black-box optimization \citep{gasnikov2022randomized}. Notably, all of these results achieve the optimal complexity derived by \citet{duchi2015optimal}. 

\subsection{On the Optimal Choices of $\{ \mu_n\}_{n=1}^\infty$ and $\{ p_n\}_{n=1}^\infty$}
In previous section, \Cref{thm:variance} connects  the perturbation stepsize sequence \( \{\mu_n\} \), the sampling distribution \( \{p_n\} \), and the variance upper bounds of our constructed unbiased estimators, which has received limited discussion in the existing literature. To control the variance term, one must ensure that $\varrho:= \sum_{n=1}^\infty\frac{(\mu_{n+1} - \mu_{n})^2}{p_n}$ (and $\varphi:=\sum_{n=1}^\infty\frac{ \mu_{n}^2}{p_n}$ for $\sfP_2$-estimator) is sufficiently small. This observation naturally raises the question: What are the optimal sequences  $\{\mu_n\}_{n=1}^\infty$ and $\{p_n\}_{n=1}^\infty$  that minimize this sum? The following theorem addresses this question:


\begin{theorem}\label{prop:lower_bound_varrho}
Let $\{\mu_n\}_{n=1}^\infty$ be a positive, decreasing sequence with   $\sum_{n=1}^\infty\mu_n<\infty$, and let $\{p_n\}_{n=1}^\infty$ be a PMF.  Denote $\mu:=\mu_1$.  Then the following statements hold:
\begin{enumerate}[label=(\alph*)]
    \item The lower bound of $\varrho$ is given by 
    $\varrho \ge \mu^2.$  
    Moreover, the equality holds if and only if
$p_n = \frac{\mu_n - \mu_{n+1}}{\mu}.$ 
    \item  The lower bound of $\varphi$ is given by 
    $\varphi\geq \Big( \sum_{n=1}^\infty \mu_n \Big)^2> \mu^2.$ 
    Moreover, the equality holds if and only if
$p_n = \frac{\mu_n }{\sum_{n=1}^\infty \mu_n }.$ 
\end{enumerate}
\end{theorem} 
% \begin{proof}
%     The lower bound and the equality condition are immediately obtained from the  Cauchy--Schwarz inequality. See \Cref{appendix:optimal-sampling} for details.
% \end{proof}

This result characterizes the choices of $\{ \mu_n\}_{n=1}^\infty$ and $\{ p_n\}_{n=1}^\infty$ that minimizes $\varrho$ (and $\varphi$ for the $\sfP_2$-estimator), leading to the variance upper bound of the form:
\begin{align*}
    \max\{\Var[\sfP_2(n,v) v] , \Var[\sfP_3(n,v) v] , \Var[\sfP_4(n,v) v] \}\leq  \calO(d\| \nabla f(x) \|^2+ d^3 \mu^2) .
\end{align*} 
Here, we can always choose $\{ \mu_n\}_{n=1}^\infty$ for the $\sfP_2$-estimator such that $\mu_1 \approx \sum_{n=2}^\infty \mu_{n}$ to nearly match the lower bound ($\approx \mu_1^2$).  

\paragraph{Sampling from the Optimal Sampling Distribution $\{ p_n\}_{n=1}^\infty$.} When the perturbation stepsize sequence $\{\mu_n\}_{n=1}^\infty$ is given, sampling the corresponding optimal distribution $p_n = \frac{\mu_n - \mu_{n+1}}{\mu_1}$ could be difficult; in most of cases, $ \{ p_n\}_{n=1}^\infty$ cannot be a ready-to-use distribution naively supported by existing software. Fortunately, we can do it conversely: given an arbitrary PMF $\{ p_n\}_{n=1}^\infty$, the perturbation stepsize takes the form
\begin{align*}
    \mu_n = \mu_1 \PP(N \geq n), \quad \text{ where } N \sim \{ p_n\}_{n=1}^\infty,
\end{align*}
providing a practical way to implement the unbiased zeroth-order gradient estimator.  
To illustrate this point, we provide two concrete examples.  


\begin{example}[Geometric $\sfP_k$-Estimators] \label{example:geometric}
We consider the geometric distribution $n \sim \operatorname{Geom}(c)$ ($c\in (0,1)$). Then    $p_n = (1-c)\,c^{\,n-1}$ for all $n\in \mathbb{N}$.  We define \(\mu_n\) by the recursion $\mu_n - \mu_{n+1} = \mu_1\,p_n = \mu_1(1-c)\,c^{\,n-1}.$ 
Summing this relation leads to the closed-form solution $$\mu_n = \mu_1\,c^{\,n-1}$$ and the optimal value $\varrho = \mu_1^2$.
This construction recovers the geometric sampling scheme used by \citet{chen2020unbiased}. We call the $\sfP_k$-estimator constructed on the geometric distribution as the \textit{geometric $\sfP_k$-estimator}. It is easy to verify that the corresponding  $\varphi$ is given as  
$\varphi = \sum_{n=1}^\infty \frac{\mu_n^2}{p_n}=\frac{\mu_1^2}{(1-c)^2}.$
\end{example}

\begin{example}[Zipf's $\sfP_k$-Estimators]\label{example:zipf}
We consider the Zipf's  distribution $n \sim \operatorname{Zipf}(s)$ ($s > 1$). Then    $p_n =  \frac{1}{\zeta(s)} \frac{1}{n^s}$ for all $n\in \mathbb{N}$, where $\zeta$ is the Riemannian zeta function defined as $\zeta(s) = \sum_{n=1}^\infty \frac{1}{n^s}.$  We define \(\mu_n\) by the recursion $\mu_n - \mu_{n+1} = \mu_1\,p_n = \mu_1 \frac{1}{\zeta(s)} \frac{1}{n^s}.$
Summing this relation leads to the closed-form solution $$\mu_n = \mu_1 \left[ 1- \frac{\sum_{j=1}^{n-1} \frac{1}{j^s}}{\zeta(s)}\right].$$
This construction also leads to the optimal value $\varrho = \mu_1^2$. When estimating the upper bound of $\varphi$, we  additionally assume $s>3$. In this case, we have
$\varphi = \sum_{n=1}^\infty \frac{\mu_n^2}{p_n}  \le\;\frac{\zeta(s-2)}{(s-1)^2\,\zeta(s)}\;\mu_1^2.$ 
The detailed calculation is put in \Cref{example:zipf-appendix}.
\end{example}   
 
In both examples, we start with a well-known easy-to-sample distribution $\{p_n\}_{n=1}^\infty$, and calculate the associated perturbation stepsize sequence $\{\mu_n\}_{n=1}^\infty$ either analytically (Geometric $\sfP_k$-estimators) or iteratively  (Zipf's $\sfP_k$-estimators). While  all estimators (i.e. $\sfP_k$-estimator with $k=2,3,4$) achieve the optimal variance in the order with $d$ and $\mu$,  these examples indicate a  key difference between the $\sfP_2$-estimator and the $\sfP_k$-estimator (for $k=3,4$): the variance bound of $\sfP_3$- and $\sfP_4$-estimator 
 is \textit{\textbf{parameter-agnostic}}; that is, once $\{p_n\}$ is specified, no additional tuning of distribution parameters is required to attain the optimal bound $\mu^2$.  This distinction highlight the practical advantages of \(\sfP_3\)- and \(\sfP_4\)-estimators. 

 


\subsection{Convergence of SGD with Unbiased Gradient Estimators}\label{sec:convergence}
In this subsection, we consider the stochastic optimization setting described in \Cref{eq:optimization-problem}, where the goal is to estimate the stochastic gradient \(\nabla f(x; \xi)\) rather than the full gradient. Under the optimal sampling distribution \(\{p_n\}_{n=1}^\infty\) and the corresponding perturbation stepsize sequence \(\{\mu_n\}_{n=1}^\infty\), the convergence upper bound of SGD follows directly from standard results for general unbiased stochastic gradient methods.   

\begin{corollary}[\citet{khaled2022better}]\label{cor:sgd-convergence}
Consider the stochastic optimization problem in \Cref{eq:optimization-problem}, and suppose that the individual loss \( f(x;\xi) \) is second-order differentiable with \( L \)-Lipschitz continuous gradient in \( x \), uniformly over \(\xi\sim\Xi\). Assume the stochastic gradient is approximated using the $\sfP_k$-estimator \( \sfP_k(n,v)\,v \) for \(k =2,3,4\). Let the SGD iteration be defined as $x_{t+1} = x_t - \eta \sfP_k(n_t, v_t)\,v_t$ where $\eta \in (0, \frac{1}{L^2d}]$ is the stepsize. Then the iterates satisfy the following convergence guarantee:
    $$\min_{0\leq t \leq T-1} \EE \|\nabla f(x_t)\|^2 \leq \calO(d^3 \mu^2  \eta + d \eta  +\frac{2}{\eta T}).$$
Consequently, choosing \(\eta = \Theta(1/\sqrt{dT})\) and $\mu = \mathcal{O}(\frac{1}{d})$ yields the optimal complexity $T = \Theta(\frac{d}{\epsilon^4})$ of having $\min_{0\leq t \leq T-1} \EE \|\nabla f(x_t)\| \leq \epsilon$. 
\end{corollary} 
 This complexity has matched the lower bound of solving a smooth non-convex optimization problem using zeroth-order gradient-based method \citep{duchi2015optimal} and cannot be further improved without adding additional assumptions.  Though we directly apply the result from \citet{khaled2022better} (which is applicable for all unbiased estimators), the zeroth-order estimation can result in an additional dependence on the dimension $d$; this dependence has been reflected in our upper bound.  
 
 
\section{Experiments}\label{sec:exp} 
To validate our theoretical results and demonstrate the effectiveness of the proposed unbiased zeroth-order gradient estimators, we conduct experiments on two settings: synthetic objectives and language model optimization. Details and hyperparameter configurations are provided in \Cref{appendix:experiments}.

 
\begin{figure}[!t]
    \centering
    \includegraphics[width=0.8\linewidth]{img/combined_dimension_comparison.png} 
    \caption{This figure presents the MSE error of four different estimators across various dimensions $d$ ranging from $16$ to $4096$. The left panel corresponds to the quadratic loss $f_{\text{reg}}$, while the right panel illustrates results for the logistic loss $f_{\text{cls}}$. Each box plot describes the distribution of the MSE error across $100$ random trials. }
    \label{fig:comparison-vs-dimension}
\end{figure}

 
\subsection{Synthetic Examples}\label{sec:synthetic}
We first evaluate our estimators on two classic loss functions \citep{james2013introduction}: the quadratic loss $f_{\text{reg}}: \RR^d \to \RR$ for linear regression and the logistic loss $f_{\text{cls}}: \RR^d \to \RR$ for binary classification. 
\begin{align*}
    f_{\text{reg}}(x)= x^\top A^\top A x, \quad f_{\text{cls}}(x) = \frac{1}{n} \sum_{i=1}^{n} \log(1 + \exp(-b_i \cdot (a_i^\top \cdot x))),
\end{align*}  
where each entry of $A \in \mathbb{R}^{d\times d}$ is independently sampled from the uniform distribution $U[-1,1]$, each feature vector $a_i \in \mathbb{R}^d$ is sampled from the standard normal distribution $\operatorname{Normal}(0, I_d)$, and $b_i \in \{-1,1\}$ are binary labels generated based on a Bernoulli distribution with the fixed sample size $n$. The gradient of each objective function can be explicitly evaluated; we compare the performance of different zeroth-order gradient estimator using the  Mean-Square-Error (MSE), which is defined as 
\begin{align}
    \text{MSE}( \hat{\nabla}f(x)):= [\hat{\nabla}f(x) - {\nabla}f(x)]^\top  [\hat{\nabla}f(x) - {\nabla}f(x)]. \label{eq:dmse}
\end{align}   

We compare the accuracy of estimating the gradient of two loss functions among four different gradient estimators including Zipf's $\sfP_3$-estimator (\Cref{example:zipf}), two-point estimator with Gaussian or uniform random perturbations, and centralized two-point estimator with uniform perturbation
(the batch size of two-point estimators is adjusted to exactly $3$ function evaluations). For detailed hyper-parameter setting, we put in \Cref{appendix:experiments}.  Several observations can be made from the results shown in \Cref{fig:comparison-vs-dimension}. First,  comparing the same estimator across different dimensions, the MSE error for both objective functions generally increases with the dimension $d$, which is expected as higher-dimensional settings pose greater estimation challenges. Second, comparing different estimators,  the Zipf's $\sfP_3$-estimator consistently achieves lower MSE compared to others.   These results collectively demonstrate the effectiveness of our proposed estimator when estimating the gradient, especially in high-dimensional settings, which will be further validate in the next experiment. 

 
 

% \begin{figure}[t] 
%         \centering
%     \includegraphics[width=0.75\linewidth]{img/a_comparison.png}
%     \caption{Enter Caption}
%     \label{fig:comparison-a}
% \end{figure}
 

% \paragraph{Sensitivity analysis on the Zeta distribution parameter}
% Furthermore, we analyze the choice of the optimal distribution \( \{p_n\}_{n=1}^\infty \) used in the P3 and P4 estimators. In our theoretical results, we show that there is no exact optimal distribution that minimizes the variance of an unbiased zeroth-order gradient estimator, due to the constraint forming a non-compact set in the functional space. Nevertheless, we observe that within certain distribution families--specifically, the zeta distribution class--it is possible to approach the limiting variance arbitrarily closely. To illustrate this phenomenon, we examine the relationship between the parameter \( \alpha \) of the zeta distribution and the MSE of the gradient estimation. As shown in \Cref{fig:comparison-a}, the MSE (corresponding to the variance of the unbiased estimator) decreases as \( \alpha \) approaches 1, and increases as \( \alpha \) becomes larger. This highlights ................ 


 

\subsection{Language Model Optimization}\label{sec:llm}
In this section, we demonstrate the practical applicability of the unbiased gradient estimators in optimizing the deep neural network. Particularly, we apply it to the task of fine-tuning a pre-trained language model.  Using zeroth-order optimization to fine-tune the LLMs has been an active research field in recent years due to its effectiveness in saving memory \citep{malladi2023fine, zhang2024revisiting, gautam2024variance, guo2024zeroth}; it allows for fine-tuning model parameters without requiring access to the full computational graph, which can be prohibitively large for modern language models. 

\begin{figure}[H]
    \centering
    \includegraphics[width=0.85\linewidth]{img/loss_vs_pseudostep_and_cost.png}  
    \caption{Comparison of training loss during fine-tuning of OPT-1.3B on SST-2 using different zeroth-order gradient estimators. The right panel rescales iterations by the number of function evaluations. The unbiased Zipf’s $\sfP_3$-, Zipf’s $\sfP_4$-, Geometric $\sfP_3$-, and Geometric $\sfP_4$-estimators achieve faster convergence under the same number of function evaluations.}
    \label{fig:lm-optimization}
\end{figure}
 
We conducted experiments using the OPT-1.3b model \citep{zhang2022opt} for sentiment classification on the Stanford Sentiment Treebank (SST-2) dataset \citep{socher-etal-2013-recursive}. To ensure fair comparison, we maintained consistent parameters across experiments: the learning rate $\eta=10^{-4}$ and the perturbation stepsize $\mu=10^{-3}$ (corresponding to $\mu_1$ in the proposed unbiased estimators), which is taken from  \citet{malladi2023fine}'s Table 7 without additional tuning. For two-point estimators, we have adjusted the batch size to align $4$ function evaluations. Detailed experimental settings are provided in \Cref{appendix:experiments}. As shown in \Cref{fig:lm-optimization}, zeroth-order optimization using the proposed unbiased zeroth-order estimators achieved superior performance compared to other baseline methods. 

\paragraph{Direct Comparison to the Two-Point Estimator ($b=1$)} Previously, we compare our proposed methods against two-point estimators under the constraint of four function evaluations. It is also interesting to consider a direct comparison with classical two-point estimators using a batch size of $b=1$, which corresponding to two function evaluations. \Cref{fig:lm-optimization-2} presents this comparison.

\begin{figure}[H]
    \centering
    \includegraphics[width=0.85\linewidth]{img/loss_vs_pseudostep_and_cost_2.png}  
    \caption{Comparison to the two-point estimator with $b=1$ under the same setting as \Cref{fig:lm-optimization}. We also include the performance of the first-order Adam and SGD in the left panel.}
    \label{fig:lm-optimization-2}
\end{figure}

 Choosing larger batch sizes gives more accurate gradient estimates, leading to lower training loss when measured by the number of updates. However, we also observe that selecting the batch size as  $b=1$  may also present its own advantage. Therefore, choosing the batch size can be non-trivial and it requires to balance the variance of gradient estimation against the per-step cost.



\section{Conclusion}
In this work, we proposed a novel class of unbiased zeroth-order gradient estimators based on a telescoping series expansion of directional derivatives. We established new theoretical results, including a sufficient condition for the expectation representation (\Cref{prop:abs-conv}), the unbiasedness of the proposed estimators (\Cref{thm:unbiased}), a variance analysis for four specific constructions (\Cref{thm:variance}), and the characterization of the optimal sampling distribution and perturbation stepsize sequence (\Cref{prop:lower_bound_varrho}). We further demonstrated that SGD equipped with our estimators achieves optimal sample complexity and empirically outperforms existing mini-batch two-point estimators. These results provide a principled foundation for a new class of estimators in zeroth-order optimization, offering both theoretical insights and practical improvements. 
 

% \section{Limitations \& Future Work}
% While our proposed estimators eliminate bias in zeroth-order gradient estimation, they require at least three function evaluations to achieve finite variance, which may limit efficiency in settings with costly evaluations. Moreover, our analysis shows that no fixed index distribution globally minimizes variance, suggesting opportunities for adaptive or learned sampling strategies. A key open question remains: 
% \centerbox{\textbf{Conjecture:} Is it possible to design a two-point unbiased zeroth-order gradient estimator with finite variance using only function evaluations?}
% Addressing this would deepen our understanding of the bias–variance trade-off in gradient-free optimization and potentially enable more efficient algorithms.
 


 




% \begin{ack}
% Use unnumbered first level headings for the acknowledgments. All acknowledgments
% go at the end of the paper before the list of references. Moreover, you are required to declare
% funding (financial activities supporting the submitted work) and competing interests (related financial activities outside the submitted work).
% More information about this disclosure can be found at: \url{https://neurips.cc/Conferences/2025/PaperInformation/FundingDisclosure}.


% Do {\bf not} include this section in the anonymized submission, only in the final paper. You can use the \texttt{ack} environment provided in the style file to automatically hide this section in the anonymized submission.
% \end{ack}
 


{
\small
\bibliography{ref} 
\bibliographystyle{plainnat} 
% \bibliographystyle{unsrt} 
}



\newpage
% \section*{NeurIPS Paper Checklist} 


% \begin{enumerate}

% \item {\bf Claims}
%     \item[] Question: Do the main claims made in the abstract and introduction accurately reflect the paper's contributions and scope?
%     \item[] Answer: \answerYes{} % Replace by \answerYes{}, \answerNo{}, or \answerNA{}.
%     \item[] Justification: We have included the main contributions in the abstract and introduction in bold. 
%     \item[] Guidelines:
%     \begin{itemize}
%         \item The answer NA means that the abstract and introduction do not include the claims made in the paper.
%         \item The abstract and/or introduction should clearly state the claims made, including the contributions made in the paper and important assumptions and limitations. A No or NA answer to this question will not be perceived well by the reviewers. 
%         \item The claims made should match theoretical and experimental results, and reflect how much the results can be expected to generalize to other settings. 
%         \item It is fine to include aspirational goals as motivation as long as it is clear that these goals are not attained by the paper. 
%     \end{itemize}

% \item {\bf Limitations}
%     \item[] Question: Does the paper discuss the limitations of the work performed by the authors?
%     \item[] Answer: \answerYes{} % Replace by \answerYes{}, \answerNo{}, or \answerNA{}.
%     \item[] Justification: We have included a separate limitation section in the appendix. 
%     \item[] Guidelines:
%     \begin{itemize}
%         \item The answer NA means that the paper has no limitation while the answer No means that the paper has limitations, but those are not discussed in the paper. 
%         \item The authors are encouraged to create a separate "Limitations" section in their paper.
%         \item The paper should point out any strong assumptions and how robust the results are to violations of these assumptions (e.g., independence assumptions, noiseless settings, model well-specification, asymptotic approximations only holding locally). The authors should reflect on how these assumptions might be violated in practice and what the implications would be.
%         \item The authors should reflect on the scope of the claims made, e.g., if the approach was only tested on a few datasets or with a few runs. In general, empirical results often depend on implicit assumptions, which should be articulated.
%         \item The authors should reflect on the factors that influence the performance of the approach. For example, a facial recognition algorithm may perform poorly when image resolution is low or images are taken in low lighting. Or a speech-to-text system might not be used reliably to provide closed captions for online lectures because it fails to handle technical jargon.
%         \item The authors should discuss the computational efficiency of the proposed algorithms and how they scale with dataset size.
%         \item If applicable, the authors should discuss possible limitations of their approach to address problems of privacy and fairness.
%         \item While the authors might fear that complete honesty about limitations might be used by reviewers as grounds for rejection, a worse outcome might be that reviewers discover limitations that aren't acknowledged in the paper. The authors should use their best judgment and recognize that individual actions in favor of transparency play an important role in developing norms that preserve the integrity of the community. Reviewers will be specifically instructed to not penalize honesty concerning limitations.
%     \end{itemize}

% \item {\bf Theory assumptions and proofs}
%     \item[] Question: For each theoretical result, does the paper provide the full set of assumptions and a complete (and correct) proof?
%     \item[] Answer: \answerYes{} % Replace by \answerYes{}, \answerNo{}, or \answerNA{}.
%     \item[] Justification: We have include the full assumptions in the statement of each lemma and theorem. 
%     \item[] Guidelines:
%     \begin{itemize}
%         \item The answer NA means that the paper does not include theoretical results. 
%         \item All the theorems, formulas, and proofs in the paper should be numbered and cross-referenced.
%         \item All assumptions should be clearly stated or referenced in the statement of any theorems.
%         \item The proofs can either appear in the main paper or the supplemental material, but if they appear in the supplemental material, the authors are encouraged to provide a short proof sketch to provide intuition. 
%         \item Inversely, any informal proof provided in the core of the paper should be complemented by formal proofs provided in appendix or supplemental material.
%         \item Theorems and Lemmas that the proof relies upon should be properly referenced. 
%     \end{itemize}

%     \item {\bf Experimental result reproducibility}
%     \item[] Question: Does the paper fully disclose all the information needed to reproduce the main experimental results of the paper to the extent that it affects the main claims and/or conclusions of the paper (regardless of whether the code and data are provided or not)?
%     \item[] Answer: \answerYes{} % Replace by \answerYes{}, \answerNo{}, or \answerNA{}.
%     \item[] Justification: We have included the full source codes and reproduction instructions. For synthetic experiment, we additionally include a separate self-contained Jupyter notebook.
%     \item[] Guidelines:
%     \begin{itemize}
%         \item The answer NA means that the paper does not include experiments.
%         \item If the paper includes experiments, a No answer to this question will not be perceived well by the reviewers: Making the paper reproducible is important, regardless of whether the code and data are provided or not.
%         \item If the contribution is a dataset and/or model, the authors should describe the steps taken to make their results reproducible or verifiable. 
%         \item Depending on the contribution, reproducibility can be accomplished in various ways. For example, if the contribution is a novel architecture, describing the architecture fully might suffice, or if the contribution is a specific model and empirical evaluation, it may be necessary to either make it possible for others to replicate the model with the same dataset, or provide access to the model. In general. releasing code and data is often one good way to accomplish this, but reproducibility can also be provided via detailed instructions for how to replicate the results, access to a hosted model (e.g., in the case of a large language model), releasing of a model checkpoint, or other means that are appropriate to the research performed.
%         \item While NeurIPS does not require releasing code, the conference does require all submissions to provide some reasonable avenue for reproducibility, which may depend on the nature of the contribution. For example
%         \begin{enumerate}
%             \item If the contribution is primarily a new algorithm, the paper should make it clear how to reproduce that algorithm.
%             \item If the contribution is primarily a new model architecture, the paper should describe the architecture clearly and fully.
%             \item If the contribution is a new model (e.g., a large language model), then there should either be a way to access this model for reproducing the results or a way to reproduce the model (e.g., with an open-source dataset or instructions for how to construct the dataset).
%             \item We recognize that reproducibility may be tricky in some cases, in which case authors are welcome to describe the particular way they provide for reproducibility. In the case of closed-source models, it may be that access to the model is limited in some way (e.g., to registered users), but it should be possible for other researchers to have some path to reproducing or verifying the results.
%         \end{enumerate}
%     \end{itemize}


% \item {\bf Open access to data and code}
%     \item[] Question: Does the paper provide open access to the data and code, with sufficient instructions to faithfully reproduce the main experimental results, as described in supplemental material?
%     \item[] Answer: \answerYes{} % Replace by \answerYes{}, \answerNo{}, or \answerNA{}.
%     \item[] Justification:  We have included the full source codes and reproduction instructions. For synthetic experiment, we additionally include a separate self-contained Jupyter notebook.
%     \item[] Guidelines:
%     \begin{itemize}
%         \item The answer NA means that paper does not include experiments requiring code.
%         \item Please see the NeurIPS code and data submission guidelines (\url{https://nips.cc/public/guides/CodeSubmissionPolicy}) for more details.
%         \item While we encourage the release of code and data, we understand that this might not be possible, so “No” is an acceptable answer. Papers cannot be rejected simply for not including code, unless this is central to the contribution (e.g., for a new open-source benchmark).
%         \item The instructions should contain the exact command and environment needed to run to reproduce the results. See the NeurIPS code and data submission guidelines (\url{https://nips.cc/public/guides/CodeSubmissionPolicy}) for more details.
%         \item The authors should provide instructions on data access and preparation, including how to access the raw data, preprocessed data, intermediate data, and generated data, etc.
%         \item The authors should provide scripts to reproduce all experimental results for the new proposed method and baselines. If only a subset of experiments are reproducible, they should state which ones are omitted from the script and why.
%         \item At submission time, to preserve anonymity, the authors should release anonymized versions (if applicable).
%         \item Providing as much information as possible in supplemental material (appended to the paper) is recommended, but including URLs to data and code is permitted.
%     \end{itemize}


% \item {\bf Experimental setting/details}
%     \item[] Question: Does the paper specify all the training and test details (e.g., data splits, hyperparameters, how they were chosen, type of optimizer, etc.) necessary to understand the results?
%     \item[] Answer:  \answerYes{} % Replace by \answerYes{}, \answerNo{}, or \answerNA{}.
%     \item[] Justification: We have included the these details in the experiment section and the appendix. 
%     \item[] Guidelines:
%     \begin{itemize}
%         \item The answer NA means that the paper does not include experiments.
%         \item The experimental setting should be presented in the core of the paper to a level of detail that is necessary to appreciate the results and make sense of them.
%         \item The full details can be provided either with the code, in appendix, or as supplemental material.
%     \end{itemize}

% \item {\bf Experiment statistical significance}
%     \item[] Question: Does the paper report error bars suitably and correctly defined or other appropriate information about the statistical significance of the experiments?
%     \item[] Answer: \answerYes{} % Replace by \answerYes{}, \answerNo{}, or \answerNA{}.
%     \item[] Justification: In the synthetic experiment, we ran 100 independent trails to obtain the estimation. The median and the $\pm 25\%$ percentiles are reported in the boxplot. 
%     \item[] Guidelines:
%     \begin{itemize}
%         \item The answer NA means that the paper does not include experiments.
%         \item The authors should answer "Yes" if the results are accompanied by error bars, confidence intervals, or statistical significance tests, at least for the experiments that support the main claims of the paper.
%         \item The factors of variability that the error bars are capturing should be clearly stated (for example, train/test split, initialization, random drawing of some parameter, or overall run with given experimental conditions).
%         \item The method for calculating the error bars should be explained (closed form formula, call to a library function, bootstrap, etc.)
%         \item The assumptions made should be given (e.g., Normally distributed errors).
%         \item It should be clear whether the error bar is the standard deviation or the standard error of the mean.
%         \item It is OK to report 1-sigma error bars, but one should state it. The authors should preferably report a 2-sigma error bar than state that they have a 96\% CI, if the hypothesis of Normality of errors is not verified.
%         \item For asymmetric distributions, the authors should be careful not to show in tables or figures symmetric error bars that would yield results that are out of range (e.g. negative error rates).
%         \item If error bars are reported in tables or plots, The authors should explain in the text how they were calculated and reference the corresponding figures or tables in the text.
%     \end{itemize}

% \item {\bf Experiments compute resources}
%     \item[] Question: For each experiment, does the paper provide sufficient information on the computer resources (type of compute workers, memory, time of execution) needed to reproduce the experiments?
%     \item[] Answer: \answerYes{} % Replace by \answerYes{}, \answerNo{}, or \answerNA{}.
%     \item[] Justification: We have provided the compute resources information in the appendix.
%     \item[] Guidelines:
%     \begin{itemize}
%         \item The answer NA means that the paper does not include experiments.
%         \item The paper should indicate the type of compute workers CPU or GPU, internal cluster, or cloud provider, including relevant memory and storage.
%         \item The paper should provide the amount of compute required for each of the individual experimental runs as well as estimate the total compute. 
%         \item The paper should disclose whether the full research project required more compute than the experiments reported in the paper (e.g., preliminary or failed experiments that didn't make it into the paper). 
%     \end{itemize}
    
% \item {\bf Code of ethics}
%     \item[] Question: Does the research conducted in the paper conform, in every respect, with the NeurIPS Code of Ethics \url{https://neurips.cc/public/EthicsGuidelines}?
%     \item[] Answer: \answerYes{} % Replace by \answerYes{}, \answerNo{}, or \answerNA{}.
%     \item[] Justification: We exactly follow the NeurIPS Code of Ethics.
%     \item[] Guidelines:
%     \begin{itemize}
%         \item The answer NA means that the authors have not reviewed the NeurIPS Code of Ethics.
%         \item If the authors answer No, they should explain the special circumstances that require a deviation from the Code of Ethics.
%         \item The authors should make sure to preserve anonymity (e.g., if there is a special consideration due to laws or regulations in their jurisdiction).
%     \end{itemize}


% \item {\bf Broader impacts}
%     \item[] Question: Does the paper discuss both potential positive societal impacts and negative societal impacts of the work performed?
%     \item[] Answer: \answerYes{} % Replace by \answerYes{}, \answerNo{}, or \answerNA{}.
%     \item[] Justification: We have included this discussion in the appendix.
%     \item[] Guidelines:
%     \begin{itemize}
%         \item The answer NA means that there is no societal impact of the work performed.
%         \item If the authors answer NA or No, they should explain why their work has no societal impact or why the paper does not address societal impact.
%         \item Examples of negative societal impacts include potential malicious or unintended uses (e.g., disinformation, generating fake profiles, surveillance), fairness considerations (e.g., deployment of technologies that could make decisions that unfairly impact specific groups), privacy considerations, and security considerations.
%         \item The conference expects that many papers will be foundational research and not tied to particular applications, let alone deployments. However, if there is a direct path to any negative applications, the authors should point it out. For example, it is legitimate to point out that an improvement in the quality of generative models could be used to generate deepfakes for disinformation. On the other hand, it is not needed to point out that a generic algorithm for optimizing neural networks could enable people to train models that generate Deepfakes faster.
%         \item The authors should consider possible harms that could arise when the technology is being used as intended and functioning correctly, harms that could arise when the technology is being used as intended but gives incorrect results, and harms following from (intentional or unintentional) misuse of the technology.
%         \item If there are negative societal impacts, the authors could also discuss possible mitigation strategies (e.g., gated release of models, providing defenses in addition to attacks, mechanisms for monitoring misuse, mechanisms to monitor how a system learns from feedback over time, improving the efficiency and accessibility of ML).
%     \end{itemize}
    
% \item {\bf Safeguards}
%     \item[] Question: Does the paper describe safeguards that have been put in place for responsible release of data or models that have a high risk for misuse (e.g., pretrained language models, image generators, or scraped datasets)?
%     \item[] Answer: \answerNA{} % Replace by \answerYes{}, \answerNo{}, or \answerNA{}.
%     \item[] Justification: The paper poses no such risks.
%     \item[] Guidelines:
%     \begin{itemize}
%         \item The answer NA means that the paper poses no such risks.
%         \item Released models that have a high risk for misuse or dual-use should be released with necessary safeguards to allow for controlled use of the model, for example by requiring that users adhere to usage guidelines or restrictions to access the model or implementing safety filters. 
%         \item Datasets that have been scraped from the Internet could pose safety risks. The authors should describe how they avoided releasing unsafe images.
%         \item We recognize that providing effective safeguards is challenging, and many papers do not require this, but we encourage authors to take this into account and make a best faith effort.
%     \end{itemize}

% \item {\bf Licenses for existing assets}
%     \item[] Question: Are the creators or original owners of assets (e.g., code, data, models), used in the paper, properly credited and are the license and terms of use explicitly mentioned and properly respected?
%     \item[] Answer: \answerYes{} % Replace by \answerYes{}, \answerNo{}, or \answerNA{}.
%     \item[] Justification: We have clearly cited the related papers of the code, data, and models. 
%     \item[] Guidelines:
%     \begin{itemize}
%         \item The answer NA means that the paper does not use existing assets.
%         \item The authors should cite the original paper that produced the code package or dataset.
%         \item The authors should state which version of the asset is used and, if possible, include a URL.
%         \item The name of the license (e.g., CC-BY 4.0) should be included for each asset.
%         \item For scraped data from a particular source (e.g., website), the copyright and terms of service of that source should be provided.
%         \item If assets are released, the license, copyright information, and terms of use in the package should be provided. For popular datasets, \url{paperswithcode.com/datasets} has curated licenses for some datasets. Their licensing guide can help determine the license of a dataset.
%         \item For existing datasets that are re-packaged, both the original license and the license of the derived asset (if it has changed) should be provided.
%         \item If this information is not available online, the authors are encouraged to reach out to the asset's creators.
%     \end{itemize}

% \item {\bf New assets}
%     \item[] Question: Are new assets introduced in the paper well documented and is the documentation provided alongside the assets?
%     \item[] Answer: \answerYes{} % Replace by \answerYes{}, \answerNo{}, or \answerNA{}.
%     \item[] Justification: All codes come with self-contained names and comments. The file, README.md, is included as the documentation.
%     \item[] Guidelines:
%     \begin{itemize}
%         \item The answer NA means that the paper does not release new assets.
%         \item Researchers should communicate the details of the dataset/code/model as part of their submissions via structured templates. This includes details about training, license, limitations, etc. 
%         \item The paper should discuss whether and how consent was obtained from people whose asset is used.
%         \item At submission time, remember to anonymize your assets (if applicable). You can either create an anonymized URL or include an anonymized zip file.
%     \end{itemize}

% \item {\bf Crowdsourcing and research with human subjects}
%     \item[] Question: For crowdsourcing experiments and research with human subjects, does the paper include the full text of instructions given to participants and screenshots, if applicable, as well as details about compensation (if any)? 
%     \item[] Answer: \answerNA{} % Replace by \answerYes{}, \answerNo{}, or \answerNA{}.
%     \item[] Justification: The paper does not involve crowdsourcing nor research with human subjects.
%     \item[] Guidelines:
%     \begin{itemize}
%         \item The answer NA means that the paper does not involve crowdsourcing nor research with human subjects.
%         \item Including this information in the supplemental material is fine, but if the main contribution of the paper involves human subjects, then as much detail as possible should be included in the main paper. 
%         \item According to the NeurIPS Code of Ethics, workers involved in data collection, curation, or other labor should be paid at least the minimum wage in the country of the data collector. 
%     \end{itemize}

% \item {\bf Institutional review board (IRB) approvals or equivalent for research with human subjects}
%     \item[] Question: Does the paper describe potential risks incurred by study participants, whether such risks were disclosed to the subjects, and whether Institutional Review Board (IRB) approvals (or an equivalent approval/review based on the requirements of your country or institution) were obtained?
%     \item[] Answer: \answerNA{} % Replace by \answerYes{}, \answerNo{}, or \answerNA{}.
%     \item[] Justification: The paper does not involve crowdsourcing nor research with human subjects.
%     \item[] Guidelines:
%     \begin{itemize}
%         \item The answer NA means that the paper does not involve crowdsourcing nor research with human subjects.
%         \item Depending on the country in which research is conducted, IRB approval (or equivalent) may be required for any human subjects research. If you obtained IRB approval, you should clearly state this in the paper. 
%         \item We recognize that the procedures for this may vary significantly between institutions and locations, and we expect authors to adhere to the NeurIPS Code of Ethics and the guidelines for their institution. 
%         \item For initial submissions, do not include any information that would break anonymity (if applicable), such as the institution conducting the review.
%     \end{itemize}

% \item {\bf Declaration of LLM usage}
%     \item[] Question: Does the paper describe the usage of LLMs if it is an important, original, or non-standard component of the core methods in this research? Note that if the LLM is used only for writing, editing, or formatting purposes and does not impact the core methodology, scientific rigorousness, or originality of the research, declaration is not required.
%     %this research? 
%     \item[] Answer: \answerNA{} % Replace by \answerYes{}, \answerNo{}, or \answerNA{}.
%     \item[] Justification: The core method development in this research does not involve LLMs as any important, original, or non-standard components.
%     \item[] Guidelines:
%     \begin{itemize}
%         \item The answer NA means that the core method development in this research does not involve LLMs as any important, original, or non-standard components.
%         \item Please refer to our LLM policy (\url{https://neurips.cc/Conferences/2025/LLM}) for what should or should not be described.
%     \end{itemize}

% \end{enumerate}



%%%%%%%%%%%%%%%%%%%%%%%%%%%%%%%%%%%%%%%%%%%%%%%%%%%%%%%%%%%%
\newpage
\appendix

\addcontentsline{toc}{section}{Appendix} % Add the appendix text to the document TOC
\part{Appendix} % Start the appendix part
\parttoc % Insert the appendix TOC
% https://math.stackexchange.com/questions/3037907/show-that-sum-limits-n-1-infty-left-fraca-n1a-n-1-right-converg
\allowdisplaybreaks
\section{Additional Backgrounds}
\subsection{Gradient Estimators in Zeroth-Order Optimization}\label{appendix:gradient-estimators}
\paragraph{One-Point Zeroth-Order Estimator} One-point estimators represent the simplest class, needing only a single function query per estimate. This construction makes them suitable when queries are costly or limited, like in online settings \citep{flaxman2005online}. A common form, motivated by Gaussian smoothing \citep{nesterov2017random}, is $$(\text{Single-Point}) \qquad \hat{\nabla}_{\text{sgl}} f(x) = \frac{1}{\mu} f(x + \mu v) v,$$ 
where $v$ is often drawn from $\operatorname{Normal}(0, I_d)$. While the expectation $\mathbb{E}[\frac{1}{\mu} f(x + \mu v) v]$ approximates the gradient of the smoothed function $\nabla_x \left[\EE_{v \sim \operatorname{Normal}(0, I_d)} 
 f(x+\mu v)\right]$, the estimator is biased regarding the true gradient $\nabla f(x)$. This bias diminishes as the smoothing parameter $\mu\to 0$ \citep{berahas2022theoretical}. However, these estimators suffer from high variance, potentially scaling with $d^2$ and exploding as $\mu \to 0$ \citep{flaxman2005online}.   

\paragraph{Two-Point Zeroth-Order Estimator} Two-point estimators improve on one-point methods by using two function evaluations, often via a finite difference along a random direction \citep{shamir2017optimal}. The standard difference form is 
\begin{align*}
    (\text{Two-Side})\qquad &\hat{\nabla}_{\text{2-side}} f(x) = \frac{f(x + \mu v)  - f(x-\mu v)}{2\mu} v, \\
    (\text{One-Side})\qquad &\hat{\nabla}_{\text{1-side}} f(x) =  \frac{f(x + \mu v)  - f(x )}{\mu} v,
\end{align*} 
requiring two queries \citep{shamir2017optimal,nesterov2017random}. This construction approximates the directional derivative \citep{chen2020unbiased}. Their expectation exactly matches the gradient of a smoothed function $\nabla_x f_\mu(x)$ \citep{nesterov2017random} and maintains a $\calO(\mu)$-level error \citep{ma2025revisiting}. Variance is significantly lower than one-point methods, often scaling linearly with dimension $d$ \citep{duchi2015optimal, berahas2022theoretical}. 

\paragraph{Multiple-Point Zeroth-Order Estimator} Multiple-point estimators use more than two function evaluations to further enhance gradient estimate quality. Common strategies include Finite-Difference method \citep{dai2015convergence} or mini-batch averaging \citep{duchi2015optimal}. The finite-difference method approximates the gradient using finite differences along each standard basis direction:
$$(\text{Finite-Difference})\qquad \hat{\nabla}_{\text{fin-diff}} f(x) = \sum_{i=1}^d \frac{f(x + \mu e_i)  - f(x-\mu e_i)}{2\mu} e_i, $$
requiring $2d$ queries, where $e_i$ is the $i$-th coordinate vector. Mini-batch averaging reduces variance by averaging $b$ independent two-point estimates: 
$$(\text{Mini-Batch})\qquad \hat{\nabla}_{\text{batch}} f(x) = \sum_{i=1}^b \frac{f(x + \mu v_i)  - f(x-\mu v_i)}{2\mu} v_i, $$
The finite-difference method offers low intrinsic variance but high query cost. Mini-batching reduces base estimator variance by $1/b$ at a cost of $b$ or $2b$ queries. These multi-point approaches can be combined arbitrary directional derivative estimators. 


\subsection{Zeroth-Order SGD and Its Variants} 

\paragraph{Vanilla Zeroth-Order SGD} The convergence of SGD has been extensively studied under various settings. \citet{ghadimi2013stochastic} established complexity results for computing approximate solutions using first-order and zeroth-order (gradient-free) information with Gaussian smoothing. For smooth convex objective functions, \citet{duchi2015optimal} obtained the optimal convergence upper bound for SGD under the zeroth-order optimization (ZOO) setting. \citet{nesterov2017random} provided the optimal convergence upper bound for Gaussian smoothing. In the realm of nonconvex optimization, \citet{ji2019improved} proposed two new zeroth-order variance-reduced optimization algorithms, ZO-SVRG-Coord-Rand and ZO-SPIDER-Coord, and provided improved analysis for the existing ZO-SVRG-Coord algorithm. These methods achieved better convergence rates and function query complexities than previous approaches. \citet{berahas2022theoretical} derived convergence analyses for finite differences, linear interpolation, Gaussian smoothing, and uniform sphere smoothing methods. Recent studies have focused on non-smooth settings. \citet{davis2022gradient} and \citet{zhang2020complexity} established the sample complexity for Lipschitz functions without assuming smoothness. \citet{lin2022gradient} derived the complexity upper bound of SGD while noting a $\sqrt{d}$ scale compared to the smooth setting. Notably, \citet{rando2023optimal} and \citet{kornowski2024algorithm} revealed that by applying certain techniques, the non-smooth case is not inherently more challenging than the smooth case. A potential direction for extending this line of research is to explore the intersection between zeroth-order SGD and random reshuffling \citep{ma2020understanding,mishchenko2020random}, minimax optimization \citep{chen2022accelerated}, or dependent data \citep{ma2022data}.
 

\paragraph{Variance-Reduced Zeroth-Order SGD}  
A key bottleneck in vanilla zeroth-order SGD is the high variance of gradient estimators, which arises from both stochastic data sampling and the inherent randomness in the gradient estimation process. This high variance necessitates small stepsizes, leading to slow convergence \citep{liu2020primer}. To address the variance from stochastic data sampling, variance reduction techniques—originally developed for first-order methods \citep{fang2018spider,defazio2014saga,johnson2013accelerating,nguyen2017sarah}, have been adapted to the zeroth-order setting. Algorithms such as ZO-SVRG \citep{liu2018zeroth,huang2019faster,gu2021black}, ZO-SVRG/SPIDER-Coord \citep{ji2019improved}, and ZO-SPIDER/SARAH \citep{fang2018spider,ji2019improved,chen2023faster} leverage epoch-based updates with variance-reducing correction terms or recursive estimator refinements. These methods significantly improve convergence by reducing the iteration complexity.

\paragraph{Memory-Efficient Zeroth-Order SGD}  
Standard SGD typically requires storing all intermediate gradient across layers to enable chain-rule-based backpropagation, which incurs substantial memory overhead, especially when training large models. To alleviate this, MeZO \citep{malladi2023fine} introduces a memory-efficient approach wherein it suffices to store the random seed used to generate the perturbation vector for each layer, dramatically reducing memory consumption. This principle motivates algorithms such as Addax \citep{li2024addax}, ElasticZO \citep{sugiura2025elasticzo}, and ZO2 \citep{wang2025zo2}, along with related benchmarking efforts \citep{zhang2024revisiting,gautam2024variance,wang2024simultaneous,guo2024zeroth}. Additional strategies exploit sparsity to further reduce memory usage; notable examples include ZORO \citep{cai2022zeroth}, the Extreme Sparsity framework \citep{guo2024zeroth}, and the One-Bit method \citep{cai2022one}.
 
\subsection{Discussions on the Forward Auto-Differentiation (AD) Approach}
The forward gradient \citep{griewank2008evaluating} provides the exact directional derivative (with exactly zero bias), while the zeroth-order approach offers only an approximation of the derivative gradient. As a result, the zeroth-order approximation inherently introduces additional variance (even if it can be unbiased). As pointed out by \citet{zhang2024revisiting}, this makes the Forward AD method theoretically better in terms of estimator quality. However, there still multiple scenarios where  the zeroth-order method is preferable.  
\begin{itemize}
    \item \textit{Implementation Difficulty}: The practical implementation of Forward AD heavily relies on the availability of JVP (a.k.a. the Jacobian-Vector Product). A naive implementation will not reduce the memory usage and potentially increase the computation cost.  

    \item \textit{Memory usage}:  Forward AD can be memory-efficient when implemented properly.  However, it still presents a higher memory usage  than the zeroth-order optimization.  Therefore, for the edge device or other extreme cases where the memory cost is sensitive, we may still prefer the zeroth-order approach.
\end{itemize}
We also note that  zeroth-order optimization is clearly advantageous in black-box settings where the forward gradient is not available. Therefore, the forward auto-differentiation and zeroth-order approaches are not mutually exclusive, but complementary, depending on the feasibility and the device memory.

\subsection{Discussions on the Difference Between Our Results and \citet{chen2020unbiased}}
Although the telescoping structure is the same as the one used in \citet{chen2020unbiased} as we have commented in the introduction section, we have developed more results to the unbiased zeroth-order gradient estimator beyond this telescoping structure:
\begin{enumerate}
    \item \textit{Identify when we can have an unbiased gradient estimator}: We identify the condition under which the telescoping structure admits a valid expectation representation (\Cref{prop:abs-conv}). This condition is critical for constructing unbiased estimators, but has not been established in \citet{chen2020unbiased}. 
    
    \item \textit{More general unbiased gradient estimator \& Reduce the number of function evaluations from $4$ to $1$:}  The estimator in \citet{chen2020unbiased} is a special case of our 
$\mathsf{P}_4$-estimator. Our framework extends beyond this, answering a fundamental theoretical question: \textit{What is the minimal number of function evaluations needed to construct an unbiased gradient estimator?} We improve the known answer from $4$ give by \citet{chen2020unbiased} to $1$. 
    \item \textit{Identify the necessary and sufficient condition of achieving the optimal variance:}  Moreover, one of our key focuses is identifying optimal parameter sequences 
 $\{\mu_n\}$ and $\{p_n\}$
 (\Cref{thm:variance}), rather than proposing a specific estimator.  
\end{enumerate}

 
\section{Bias Analysis}\label{appendix:bias}
\subsection{Absolute Convergence}\label{appendix:absolute-convergence}%ALL CLEAR
In this subsection, we derive a sufficient condition to guarantee the expectation representation of the telescoping series introduced in \Cref{eq:expectation}. This requires ensuring the series is absolutely convergent, a property essential for interpreting it as the expectation of a well-defined random variable. The following definitions are directly taken from \citet{Folland_Taylor2024}:
\begin{definition}[Convergent series]
    A series $\sum_{n=1}^\infty a_n$ is said to be \textit{convergent} if the sequence of partial sums $S_N=\sum_{n=1}^N a_n$ is a convergent sequence; that is, $\lim_{N\to \infty} S_N$ exists.
\end{definition}
\begin{definition}[Absolutely convergent series]
    A series $\sum_{n=1}^\infty a_n$ is said to be \textit{absolutely convergent} if the series $\sum_{n=1}^\infty |a_n|$ is convergent. 
\end{definition}

The following classical result, known as the Riemann series theorem \citep{riemann1868darstellbarkeit,spivak2008calculus}, highlights the necessity of absolute convergence when interpreting an infinite sum as a well-defined value. 
\begin{theorem}[Riemann Series Theorem]
    Let \( \sum_{n=1}^\infty a_n \) be a conditionally convergent series of real numbers (i.e., convergent but not absolutely convergent). Then, for any real number \( r \in \RR \), there exists a rearrangement \( \sigma: \mathbb{N} \to \mathbb{N} \) such that $\sum_{n=1}^\infty a_{\sigma(n)} = r.$ Moreover, there exist rearrangements such that the sum diverges to \( +\infty \), \( -\infty \), or fails to converge at all.
\end{theorem}
This theorem underscores the critical distinction between conditional and absolute convergence: a convergent series may yield different values under different summation orders. However, the definition of expectation for a random variable does not permit such ambiguity, since the outcomes have no inherent order. To guarantee a well-defined expectation, absolute convergence is required: that is, for a random variable $X$ with outcomes $\{x_n\}$ and probabilities $\{p_n\}$, it must hold that $\EE |X| = \sum_{n=1}^\infty |x_n| p_n < \infty$. 

We now recap \Cref{prop:abs-conv} describing the condition where the telescoping series is absolutely convergent, enabling a valid expectation representation.
\begin{proposition}\label{prop:abs-conv-appendix}
If the second-order continuously differentiable function $f:\RR^d \to \RR$ has $L$-Lipschitz continuous gradient and $\sum_{n=1}^\infty \mu_n < \infty$, then the series $$\sum_{n=1}^\infty p_n\left[\frac{f(x+\mu_1 v) - f(x)}{\mu_1} + \frac{1}{p_n}\left( \frac{f(x+\mu_{n+1} v) - f(x)}{\mu_{n+1}} - \frac{f(x+\mu_n v) - f(x)}{\mu_n}\right) \right]$$ is absolutely convergent and its limit is $\nabla_v f(x)$.
\end{proposition}
\begin{proof}  First, because $\sum_{n=1}^\infty |a_n+b_n| \leq \sum_{n=1}^\infty |a_n | +\sum_{n=1}^\infty |b_n |$, it suffices to prove $$\sum_{n=1}^\infty  \left( \frac{f(x+\mu_{n+1} v) - f(x)}{\mu_{n+1}} - \frac{f(x+\mu_n v) - f(x)}{\mu_n}\right)  $$ is absolutely convergent. To prove its absolute convergence, we estimate the magnitude of the difference term using Taylor’s theorem: 
    \begin{align*}
        &\left| \frac{f(x+\mu_{n+1} v)-f(x)}{\mu_{n+1}}-\frac{f(x+\mu_n v)-f(x)}{\mu_n}\right| \\
        \overset{(i)}{=} & \left| \frac{ \mu_{n+1}\nabla f(x)^\top v  + R(x) \mu_{n+1}^2}{ \mu_{n+1}} - \frac{ \mu_{n}\nabla f(x)^\top v  + R'(x) \mu_{n}^2}{ \mu_{n}}\right| \\ 
        \overset{(ii)}{\leq} & \frac{L}{2}\left|   \mu_{n+1}  +  \mu_{n} \right|
    \end{align*} 
    where (i) applies the Taylor theorem \citep{Folland_Taylor2024} with setting the integral form remainder $R(x):=\int_0^1 (1-t) \sum_{|\alpha|=2} \frac{\partial^2}{\partial x_1^{\alpha_1}\dots x_d^{\alpha_d}}   f(x+t\mu v) \ dt$, (ii) assumes the global Lipschitz continuous gradient, which results in the uniform estimate
    $R(x) \leq \frac{L}{2}$
    for all $x \in \RR^d$.  Therefore, to ensure the telescoping series is absolutely continuous, it suffices to require $\sum_{n=1}^\infty \mu_n < \infty.$  
    The limit is directly determined by the original convergent series. It concludes our proof. 
\end{proof} 

In the following two examples, we present commonly used sequences \( \{\mu_n\} \) that satisfy the condition \( \sum_{n=1}^\infty \mu_n < \infty \), ensuring the absolute convergence required in \Cref{prop:abs-conv-appendix}.

\begin{example}[Exponential Decay]
Let \( \mu_n = \alpha^n \) for some constant \( 0 < \alpha < 1 \). Then,
\[
\sum_{n=1}^\infty \mu_n = \sum_{n=1}^\infty \alpha^n = \frac{\alpha}{1 - \alpha} < \infty.
\] 
\end{example}

\begin{example}[Polynomial Decay]
Let \( \mu_n = \frac{1}{n^s} \) for some constant \( s > 1 \). Then,
\[
\sum_{n=1}^\infty \mu_n = \sum_{n=1}^\infty \frac{1}{n^s} < \infty,
\]
which is well-known as the Riemann zeta function $\zeta(s)$. 
\end{example} 

  

\subsection{Unbiasedness of Zeroth-Order Estimators in $\mathscr{P}$-Family} \label{sec:unbiasedness}  
We begin by recalling the definition of the unbiased zeroth-order gradient estimators, denoted by $\mathscr{P}:=\mathscr{P}(f, \{\mu_n\}_{n=1}^\infty, \{p_n\}_{n=1}^\infty,V)$, as given by \Cref{def:unbiased-family}. Under suitable conditions on the differentiable function $f:\RR^d \to \RR$ and the sequence $\{\mu_n\}_{n=1}^\infty$ (e.g., those provided by \Cref{prop:abs-conv}), this definition naturally yields the desired expectation representation. Moreover,  the random distributions $\{p_n\}_{n=1}^\infty$ and $V$ are independent. These conditions are sufficient to guarantee the unbiasedness stated in the following result.
\begin{theorem}[Unbiasedness]\label{thm:unbiased-appendix}
Let $\mathscr{P}:=\mathscr{P}(f, \{ \mu_n\}_{n=1}^\infty, \{ p_n\}_{n=1}^\infty,V)$ is defined as \Cref{def:unbiased-family}. Then, for any estimator $\mathsf{P}(n,v) \in \mathscr{P}$, the following hold:
\begin{enumerate}[label=(\alph*)]
    \item \(\mathbb{E}[\mathsf{P}(n,v) \mid v] = \nabla_v f(x)\); that is, \(\mathsf{P}(n,v)\) is an unbiased estimator of the directional derivative $\nabla_v f(x)$.
    \item If the random direction \(v\) is chosen independently of the sampling $n\sim \{p_n\}_{n=1}^\infty $ and satisfies \(\mathbb{E}[v\,v^\top]=I\), then
    \[
    \mathbb{E}_{n\sim \{p_n\}_{n=1}^\infty , v\sim V}\Bigl[\mathsf{P}(n,v)\,v\Bigr] = \nabla f(x),
    \]
    so that \(\mathsf{P}(n,v)\,v\) is an unbiased estimator of the full gradient.
\end{enumerate}
\end{theorem}
\begin{proof} By \Cref{def:unbiased-family}, the directional derivative  $\nabla_v f(x)$ naturally has the expectation representation. 
\begin{enumerate}[label=(\alph*)]
    \item Denote
    $$X_n:=\frac{f(x+\mu_1 v) - f(x)}{\mu_1} + \frac{1}{p_n}\left( \frac{f(x+\mu_{n+1} v) - f(x)}{\mu_{n+1}} - \frac{f(x+\mu_n v) - f(x)}{\mu_n}\right).$$
    Then by the tower property of the conditional expectation,
    \begin{align*}
        \nabla_v f(x) =\EE_{n\sim  \{ p_n\}_{n=1}^\infty}[X_n\mid v]  = \EE_{n\sim  \{ p_n\}_{n=1}^\infty}[ \EE[ P(n,v)| n] \mid v] = \EE[P(n,v)\mid v].
    \end{align*}
    It concludes the proof. 

    \item We consider the conditional expectation $\EE[\cdot|v]$. We have
    \begin{align*}
        \EE\Bigl[\mathsf{P}(n,v)\,v\Bigr] &= \EE \left[ \EE\Bigl[\mathsf{P}(n,v)\,v \Big| v \Bigr]  \right]\\ 
        &=\EE \left[ \EE\Bigl[\mathsf{P}(n,v) \Big| v \Bigr] v  \right] \\ 
        &= \EE[ \nabla_v f(x) v] \\ 
        &= \nabla f(x) 
    \end{align*}
    Therefore, we conclude that $\mathsf{P}(n,v) v$ is an unbiased estimator of the gradient $\nabla f(x)$.  
\end{enumerate}
\end{proof} 
    

\section{Variance Analysis}\label{appendix:variance} 

In this section we present the proof of \Cref{thm:variance}. Each subsection contains the proof of the corresponding item. We recap its statement below:
\begin{theorem}\label{thm:variance-appendix} 
Let $\mathscr{P}:=\mathscr{P}(f, \{ \mu_n\}_{n=1}^\infty, \{ p_n\}_{n=1}^\infty,V)$ is defined as \Cref{def:unbiased-family}. Suppose that $f:\RR^d \to \RR$ is second-order continuously differentiable and has $L$-Lipschitz continuous gradient, $\sum_{n=1}^\infty \mu_n < \infty$, and $V$ is the uniform distribution over the  sphere with the radius $\sqrt{d}$. Define 
\begin{align*}
    \mu:=\mu_1, \qquad\varrho:= \sum_{n=1}^\infty \frac{(\mu_{n+1} - \mu_{n})^2}{p_n}, \qquad \text{and} \qquad \varphi:=\sum_{n=1}^\infty  \frac{\mu_{n}^2}{p_n}.
\end{align*}
Then the following statements hold:
    \begin{enumerate}[label=(\alph*)]
        \item If there exists a point $x\in \RR^d$ such that the Hessian $\nabla^2 f(x)$ is positive definite and $f(x)\neq 0$, then the variances of the $\sfP_1$ is infinite. That is,
        $$\Var[\sfP_1(n,v) v] = + \infty.$$
        \item  The variance of $\sfP_2$-estimator $\sfP_2(n,v) v$ is given by
\begin{align*}
    \Var[\sfP_2(n,v) \, v ]  
    \leq  \Var[\sfP_4(n,v)\, v] + \frac{L^2 }{3} d^3\mu^2 + \frac{L^2}{12}  d^3  \varrho  + \frac{L^2}{3} d^3   \varphi.
\end{align*}
        \item  The variance of $\sfP_3$-estimator $\sfP_3(n,v) v$ is given by
        $$\Var[\sfP_3(n,v) v] \leq \Var[\sfP_4(n,v)\, v] + \frac{L^2 }{8} d^3\mu^2 + \frac{L^2}{8}  d^3 \varrho. $$ 
        \item The variance of $\sfP_4$-estimator $\sfP_4(n,v) v$ is given by
        $$\Var[\sfP_4(n,v) v] \leq  (d-1) \| \nabla f(x) \|^2 +  \frac{3L^2}{4} d^3 \mu^2  +  \frac{L^2 d^3}{2}  \varrho.$$
    \end{enumerate}
\end{theorem}
\begin{proof} For the item (a), we present the proof in \Cref{lemma:variance-p1}.  For arbitrary $\sfP:= \sfP(n,v)\in \mathscr{P}$, we have
\begin{align*}
    \Var[\sfP \, v ]  \overset{(i)}{=} & \EE[\sfP^2 v^\top v] - \|\nabla f(x)\|^2 \\ 
     = & d\EE[(\sfP- \sfP_4(n,v) +  \sfP_4(n,v) )^2  ] - \|\nabla f(x)\|^2 \\ 
     \overset{(ii)}{=}&d\EE[(\sfP- \sfP_4(n,v))^2  ] + d\EE[ \sfP_4(n,v) ] - \|\nabla f(x)\|^2  \\ 
     = & d\EE[(\sfP- \sfP_4(n,v))^2  ] + \Var[\sfP_4(n,v)\, v]
\end{align*}
where (i) applies the unbiasedness (\Cref{thm:unbiased}) and (ii) applies the definition of  $\mathscr{P}$ ($\sfP$ is an unbiased estimator of $\sfP_4$). Therefore, we start with $\Var[\sfP_4\, v]$, the variance of the $\sfP_4$-estimator. Then it suffices to evaluate $\EE[(\sfP_k(n,v)- \sfP_4(n,v))^2  ]$ for $k=2$ and $k=3$. Therefore, we prove the item (d) first. The detailed proof is included in \Cref{lemma:variance-p4}. Based on this result, we obtain the variance upper bound of $\sfP_2$- and $\sfP_3$-estimators in \Cref{lemma:variance-p2} and \Cref{lemma:variance-p3}, respectively.
\end{proof}

\subsection{Variance of $\mathsf{P}_1$-Estimator} %ALL CLEAR

\begin{lemma}\label{lemma:variance-p1}
    Under the same setting as \Cref{thm:variance-appendix}, if there exists a point $x\in \RR^d$ such that the Hessian $\nabla^2 f(x)$ is positive definite, then the variances of the $\sfP_1$--estimator at $x$ is infinite.  
\end{lemma}
\begin{proof}
Recall that for the random direction $\frac{1}{\sqrt{d}}v\sim \Unif(\mathbb{S}^{d-1})$ and the random variable $\mathrm{U}_4 \sim \operatorname{Uniform}\left( \{ 0, 1,2,3 \}\right)$, we have the $\sfP_1$\nobreakdash-estimator defined as
\begin{align*} 
\mathsf{P}_1(n, v)
= 
&\;\frac{f(x+\mu_1 v)\mathbb{I}_{\{\mathrm{U}_4=1\}}-f(x)\mathbb{I}_{\{\mathrm{U}_4=0\}}}{\mu_1}  \\ 
&+\frac{1}{p_n}\Biggl[
    \frac{f(x+\mu_{n+1} v)\mathbb{I}_{\{\mathrm{U}_4=2\}}-f(x)\mathbb{I}_{\{\mathrm{U}_4=0\}}}{\mu_{n+1}}
  - \frac{f(x+\mu_n v)\mathbb{I}_{\{\mathrm{U}_4=3\}}-f(x)\mathbb{I}_{\{\mathrm{U}_4=0\}}}{\mu_n}
\Biggr].
\end{align*} 
Since this estimator is unbiased (\Cref{thm:unbiased}),  
\[
\Var\bigl[\mathsf{P}_1(n,v)\,v\bigr]
= d\EE\bigl[\mathsf{P}_1(n,v)^2\bigr] - \|\nabla f(x)\|^2,
\]
so it suffices to show $\EE[\mathsf{P}_1(n,v)^2]=\infty$.  For brevity we write
\[
\sfP_1 := \mathsf{P}_1(n,v).
\]
As the Hessian $\nabla^2 f(x)$ is positive definite at the point $x$, it has $\nabla^2 f(x)\succeq\lambda_{\min}I$ for some $\lambda_{\min}>0$.  
We consider the event $\{N=n,\;\mathrm{U}_3=3\}$, one finds
\[
\sfP_1 = -\,\frac{f(x+\mu_n v)}{p_n \mu_n}
\quad\text{with probability}\quad p_n\cdot\tfrac14.
\]
Hence
\begin{align*}
    \EE_{n\sim \{ p_n\}}[\sfP_1^2] &\geq \sum_{n=1}^\infty
\Pr(N=n,\mathrm{U}_3=3)\,\Bigl(\frac{f(x+\mu_n v)}{p_n \mu_n}\Bigr)^2 \\
&\geq  \frac{1}{3} \sum_{n=1}^\infty \frac{\left[ f(x+\mu_n v)\right]^2}{p_n \mu_n^2}\\
&\geq  \frac{1}{3} \sum_{n=1}^\infty\left(  \frac{\left[  f(x)\right]^2}{p_n \mu_n^2} + \frac{\left[  f(x+\mu_n v) - f(x) \right] f(x) }{p_n \mu_n^2} \right)
\end{align*}
Without loss of generality, we assume $f(x)>0$. In the case $f(x)<0$, we apply the $L$-smoothness to obtain a lower bound instead. By the second‐order Taylor expansion \citep{spivak2008calculus} (or the strong convexity near the point $x$),
$$f(x+\mu v)
\geq f(x) + \mu\,\nabla f(x)^\top v
  + \tfrac12\,\mu^2 \lambda_{\min} \,v^\top \,v .$$
so we have
$$\frac{\left[  f(x+\mu_n v) - f(x) \right] f(x) }{p_n \mu_n^2}  \geq \frac{\nabla f(x)^\top v}{p_n \mu_n} + \frac{d \lambda_{\min } f(x)}{2p_n}.$$
As the result, we have
\begin{align*}
    \EE_{n\sim \{ p_n\},\frac{1}{\sqrt{d}}v\sim \operatorname{Uniform}(\mathbb{S}^{d-1})}[\sfP_1^2] &\geq  \frac{1}{3}\sum_{n=1}^\infty \frac{1}{p_n} \left( \frac{f(x)^2}{\mu_n^2} +  \frac{d \lambda_{\min } f(x)}{2p_n}\right).
\end{align*}
As $\{p_n\}$ is a PMF, it must diverge to infinite.  
\end{proof}


\subsection{Variance of $\mathsf{P}_4$-Estimator}%ALL CLEAR

\begin{lemma}\label{lemma:variance-p4}
    Under the same setting as \Cref{thm:variance-appendix}, the variance of $\sfP_4$-estimator $\mathsf P_4(n,v)\,v$ is upper bounded by
\begin{align*}
      \Var[\mathsf P_4(n,v)\,v  ] & \leq (d-1) \| \nabla f(x) \|^2 + \frac{3L^2 d^3 \mu_1^2}{4}  +  \frac{L^2 d^3}{2}\sum_{n=1}^\infty \frac{| \mu_{n+1} - \mu_{n}|^2}{p_n} .
\end{align*} 
\end{lemma}
\begin{proof}

Recall that
\[
\mathsf P_4(n,v)
= \frac{f(x+\mu_1v)-f(x)}{\mu_1}
+ \frac1{p_n}\Bigl[\frac{f(x+\mu_{n+1}v)-f(x)}{\mu_{n+1}}
- \frac{f(x+\mu_nv)-f(x)}{\mu_n}\Bigr].
\]
Our goal is to bound \(\Var[\mathsf P_4(n,v)\,v]\). For each \(n\), define the remainder term
\begin{align*}
    \delta_n(v)
=&\frac{f(x+\mu_n v)-f(x)-\mu_n\,\nabla f(x)^\top v}{\mu_n}, \\
 \Delta_n =& \frac{f(x+\mu_n v)-f(x) }{\mu_n} = \delta_n(v) + \nabla f(x)^\top v.
\end{align*} 
Then
\[
\mathsf P_4(n,v)
= \nabla f(x)^\top v
+ \delta_1(v)
+ \frac{\delta_{n+1}(v)-\delta_n(v)}{p_n}.
\]
Hence our vector estimator is
\[
\mathsf P_4(n,v)\,v
= vv^\top  \nabla f(x) + \bigl[\delta_1(v)+\tfrac{\delta_{n+1}(v)-\delta_n(v)}{p_n}\bigr]\,v.
\]
Since we have shown that $\mathsf P_4(n,v)\,v$ is unbiased (\Cref{thm:unbiased}) and $\EE[vv^\top]=I_d$, we have 
\begin{align}\label{eq:1}
    \EE \bigl[\delta_1(v)+\tfrac{\delta_{n+1}(v)-\delta_n(v)}{p_n}\bigr]\,v = 0.
\end{align}
The variance is given by
\begin{align*}
    \Var[\mathsf P_4(n,v)\,v  ] &= d \EE [\mathsf P_4(n,v)^2 ] - \|\nabla f(x) \|^2 \\ 
    &= (d-1)\|\nabla f(x) \|^2  + d \EE\left(\delta_1(v)
+ \frac{\delta_{n+1}(v)-\delta_n(v)}{p_n} \right)^2  \\ 
&\quad + 2 d \EE \nabla f(x)^\top v \Big(\delta_1(v) + \frac{\delta_{n+1}(v)-\delta_n(v)}{p_n}  \Big) \\
    &\overset{(i)}{=} (d-1)\|\nabla f(x) \|^2  +d \EE\left(\delta_1(v)
+ \frac{\delta_{n+1}(v)-\delta_n(v)}{p_n} \right)^2  \\  
&= (d-1)\|\nabla f(x) \|^2  + d \EE \delta_1^2(v) \\ 
& \quad + 2 d \EE \Big[\frac{\delta_1(v)}{p_n} \left( \delta_{n+1}(v)-\delta_n(v)\right)  \Big]+    \EE \Big[\frac{d}{p_n^2} \left( \delta_{n+1}(v)-\delta_n(v)\right)^2 \Big] \\ 
&\overset{(ii)}{=} (d-1)\|\nabla f(x) \|^2  + 3d\EE \delta_1^2(v)  +    \EE \Big[\frac{d}{p_n^2} \left( \delta_{n+1}(v)-\delta_n(v)\right)^2 \Big] \\ 
\end{align*}
where (i) applies the unbiasedness \Cref{eq:1} and (ii) applies the telescoping series \Cref{eq:telescoping}. Now it remains to upper bound the remainder term $\delta_n(v)$ and the  remainder difference term $\delta_{n+1}(v) - \delta_{n}(v) $.
\begin{itemize}
    \item \textbf{Bound the remainder term $\delta_n(v)$: } By $L$-smoothness of $f:\RR^d \to \RR$, we have
\begin{align} 
    f(x+\mu_n v) - f(x) &\leq  \mu_n \nabla f(x)^\top v + \frac{L \mu_n^2}{2} \|v\|^2 ,\nonumber\\ 
    &\overset{(i)}{=} \mu_n \nabla f(x)^\top v + \frac{dL \mu_n^2}{2}  , \nonumber
\end{align}
where (i) applies  $\EE vv^\top = I_d$ (this condition ensures that $\| v \|^2=\operatorname{Tr}\left(\| v \|^2\right) = \operatorname{Tr}\left( vv^\top \right) = d$).  As the result,
\begin{align}\label{eq:delta}
    \delta_n(v) \leq \frac{Ld}{2} \mu_n.
\end{align} 
or we may use the following almost-sure upper bound
\begin{align*} 
   \left[  \delta_n(v) \right]^2 \leq \frac{L^2d^2}{4} \mu_n^2.
\end{align*} 


    \item  \textbf{Bound the remainder difference term $\delta_{n+1}(v) - \delta_{n}(v) $: }  We define the remainder term as a function in $\mu$; that is,
    $$\phi(\mu):= \frac{f(x+\mu v)-f(x)-\mu\,\nabla f(x)^\top v}{\mu}.$$
    It automatically gives
    \begin{align*}
        \delta_{n+1}(v) -  \delta_{n}(v)  = \phi(\mu_{n+1}) - \phi(\mu_{n}).
    \end{align*}
    As $\phi: [\mu_{n+1}, \mu_{n}] \to \RR$ is a continuous differentiable function (by our assumption that $f:\RR^d \to \RR$ is second-order continuously differentiable), we can apply the mean-value theorem \citep{Folland_Taylor2024}: There exists $\varsigma \in [\mu_{n+1}, \mu_{n}]$ such that
    $$\phi(\mu_{n+1}) - \phi(\mu_{n}) = \phi'(\varsigma) (\mu_{n+1} - \mu_{n}).$$
    We again applying the $L$-smoothness of $f:\RR^d \to \RR$ (essentially the bounded Hessian assumption) to $\phi'(\varsigma)$, which leads to
    \begin{align}\label{eq:delta-difference}
        |\delta_{n+1}(v) -  \delta_{n}(v) | \leq \frac{Ld}{2}| \mu_{n+1} - \mu_{n}|.
    \end{align} 
\end{itemize}
As the result, we obtain the upper bound as 
\begin{align*}
      \Var[\mathsf P_4(n,v)\,v  ] &\overset{(i)}{\leq } (d-1) \| \nabla f(x) \|^2 + \frac{3L^2 d^3 \mu_1^2}{4}   + \sum_{n=1}^\infty \frac{2d}{p_n} \left(  \frac{Ld}{2} | \mu_{n+1} - \mu_{n}|\right)^2\\
      &\leq (d-1) \| \nabla f(x) \|^2 + \frac{3L^2 d^3 \mu_1^2}{4}  +  \frac{L^2 d^3}{2}\sum_{n=1}^\infty \frac{| \mu_{n+1} - \mu_{n}|^2}{p_n} ,
\end{align*} 
where (i) applies the expectation over $n\sim \{p_n\}_{n=1}^\infty$, which cancels out one $\frac{1}{p_n}$. 
\end{proof}


\subsection{Variance of $\mathsf{P}_2$-Estimator} %ALL CLEAR

\begin{lemma}\label{lemma:variance-p2}
    Under the same setting as \Cref{thm:variance-appendix}, the variance of $\sfP_2$-estimator $\mathsf P_2(n,v)\,v$ is upper bounded by
\begin{align*}
    \Var[\sfP_2(n,v) \, v ]  
    \leq  \Var[\sfP_4(n,v)\, v] + \frac{L^2}{12}  d^3\sum_{n=1}^\infty \frac{(\mu_n - \mu_{n+1})^2}{p_n} + \frac{L^2 }{3} d^3\mu_1^2 + \frac{L^2}{3} d^3 \sum_{n=1}^\infty \frac{\mu_n^2}{p_n}.
\end{align*}
\end{lemma}
\begin{proof}
Recall that $\mathsf{P}_2(n,v)$ is defined as 
\begin{align*}
    \mathsf{P}_2(n, v)
= 
&\;\frac{f(x+\mu_1 v)-f(x)}{\mu_1}\,\mathbb{I}_{\{\mathrm{U}_3=0\}}  \\ 
&\;+\;\frac{1}{p_n}\Biggl[
    \frac{f(x+\mu_{n+1} v)-f(x)}{\mu_{n+1}}\mathbb{I}_{\{\mathrm{U}_3=1\}}
  - \frac{f(x+\mu_n v)-f(x)}{\mu_n}\mathbb{I}_{\{\mathrm{U}_3=2\}}
\Biggr],
\end{align*} 
where $\mathrm{U}_3\sim  \operatorname{Uniform}\left( \{ 0, 1,2 \}\right)$ is a selection variable. Then 
\begin{align*}
    &\EE[ (\mathsf{P}_2(n, v) - \mathsf{P}_4(n, v))^2\mid n,v] \\
    = &\PP(\mathrm{U}_3=0) \left[  \frac{1}{p_n}\Biggl[ \frac{f(x+\mu_{n+1} v)-f(x)}{\mu_{n+1}} - \frac{f(x+\mu_n v)-f(x)}{\mu_n} \Biggr] \right]^2  \\ 
    &+\PP(\mathrm{U}_3=1) \left[\frac{f(x+\mu_1 v)-f(x)}{\mu_1}  +    \frac{1}{p_n}\Biggl[  - \frac{f(x+\mu_n v)-f(x)}{\mu_n} \Biggr] \right]^2  \\ 
      &+\PP(\mathrm{U}_3=2) \left[\frac{f(x+\mu_1 v)-f(x)}{\mu_1}  +   \frac{1}{p_n}\Biggl[ \frac{f(x+\mu_{n+1} v)-f(x)}{\mu_{n+1}} \Biggr]  \right]^2  \\  
    \overset{(i)}{\leq} & \frac{1}{3} \frac{1}{p_n^2} \left( \delta_{n+1}(v) - \delta_n(v)\right)^2 + \frac{2}{3} \left( [\delta_1(v)]^2 + \frac{1}{p_n^2} [\delta_n(v)]^2 \right)+ \frac{2}{3} \left( [\delta_1(v)]^2 + \frac{1}{p_n^2} [\delta_{n+1}(v)]^2 \right)\\
   = &  \frac{4}{3}[\delta_1(v)]^2  + \frac{1}{3} \frac{1}{p_n^2} \left( \delta_{n+1}(v) - \delta_n(v)\right)^2 + \frac{2}{3} \frac{1}{p_n^2} [\delta_{n}(v)]^2 +  \frac{2}{3} \frac{1}{p_n^2} [\delta_{n+1}(v)]^2 ,
    \end{align*}
    where (i) applies $(a+b)^2 \leq 2a^2 + 2b^2$ and $\delta_n(v):= \frac{f(x+\mu_n v) - f(x) -\mu_n \nabla f(x)^\top v}{\mu_n}$. As we have bounded this term in \Cref{lemma:variance-p4}, we have
    \begin{align*}
        &\quad \EE[ (\mathsf{P}_2(n, v) - \mathsf{P}_4(n, v))^2\mid  v] \\
        &= \sum_{n=1}^\infty p_n [ (\mathsf{P}_2(n, v) - \mathsf{P}_4(n, v))^2\mid  n, v] \\ 
        &\leq \sum_{n=1}^\infty p_n \left[ \frac{L^2 d^2}{3}\mu_1^2 + \frac{L^2 d^2}{12}\frac{(\mu_n-\mu_{n+1})^2}{p_n^2} + \frac{L^2 d^2}{3 p_n^2} \mu_{n}^2\right]  \\
        &\leq \frac{L^2}{12}  d^2\sum_{n=1}^\infty \frac{(\mu_n - \mu_{n+1})^2}{p_n} + \frac{L^2 }{3} d^2\mu_1^2 + \frac{L^2}{3} d^2 \sum_{n=1}^\infty \frac{\mu_n^2}{p_n}.
    \end{align*}
    Therefore, we finally obtain
\begin{align*}
    \Var[\sfP_2(n,v) \, v ]  = & d\EE[(\sfP_2(n,v)- \sfP_4(n,v))^2  ] + \Var[\sfP_4(n,v)\, v] \\ 
    \leq & \Var[\sfP_4(n,v)\, v] + \frac{L^2}{12}  d^3\sum_{n=1}^\infty \frac{(\mu_n - \mu_{n+1})^2}{p_n} + \frac{L^2 }{3} d^3\mu_1^2 + \frac{L^2}{3} d^3 \sum_{n=1}^\infty \frac{\mu_n^2}{p_n}.
\end{align*}
It completes the proof. 
\end{proof}

 




\subsection{Variance of $\mathsf{P}_3$-Estimator}%ALLCLEAR
 
\begin{lemma}\label{lemma:variance-p3}
    Under the same setting as \Cref{thm:variance-appendix}, the variance of $\sfP_3$-estimator $\mathsf P_3(n,v)\,v$ is upper bounded by
\begin{align*}
        \Var[\sfP_3(n,v) \, v ]  
 \leq \Var[\sfP_4(n,v)\, v] + \frac{L^2}{8}  d^3\sum_{n=1}^\infty \frac{(\mu_n - \mu_{n+1})^2}{p_n} + \frac{L^2 }{8} d^3\mu_1^2.
\end{align*} 
\end{lemma}
\begin{proof}
Recall that $\mathsf{P}_3(n,v)$ is defined as 
\begin{align*}
&\mathsf{P}_3(n, v)\\
= &\frac{f(x+\mu_1 v)-f(x)}{\mu_1}\mathrm{U}_2 + \frac{1}{p_n}\Biggl[ \frac{f(x+\mu_{n+1} v)-f(x)}{\mu_{n+1}} - \frac{f(x+\mu_n v)-f(x)}{\mu_n} \Biggr](1-\mathrm{U}_2),
\end{align*}
where $\mathrm{U}_2\sim  \operatorname{Uniform}\left( \{ 0, 1 \}\right)$ is a selection variable. Then 
\begin{align*}
    &\EE[ (\mathsf{P}_3(n, v) - \mathsf{P}_4(n, v))^2\mid n,v] \\
    = &\PP(\mathrm{U}_2=0) \left[ \frac{1}{p_n}\Biggl[ \frac{f(x+\mu_{n+1} v)-f(x)}{\mu_{n+1}} - \frac{f(x+\mu_n v)-f(x)}{\mu_n} \Biggr] - \mathsf{P}_4(n, v)\right]^2\\ 
    &+ \PP(\mathrm{U}_2=1)\left[\frac{f(x+\mu_1 v)-f(x)}{\mu_1}  - \mathsf{P}_4(n, v)\right]^2 \\ 
    = &\frac{1}{2} \left[ \frac{1}{p_n}\Biggl[ \frac{f(x+\mu_{n+1} v)-f(x)}{\mu_{n+1}} - \frac{f(x+\mu_n v)-f(x)}{\mu_n} \Biggr] \right]^2 + \frac{1}{2}\left[\frac{f(x+\mu_1 v)-f(x)}{\mu_1}  \right]^2 .\end{align*}
    As we have bounded this term in \Cref{lemma:variance-p4}, we have
    \begin{align*}
        \EE[ (\mathsf{P}_3(n, v) - \mathsf{P}_4(n, v))^2\mid  v] &= \sum_{n=1}^\infty p_n [ (\mathsf{P}_3(n, v) - \mathsf{P}_4(n, v))^2\mid  n, v] \\ 
        &\leq \sum_{n=1}^\infty p_n \left[  \frac{1}{ p_n^2} \frac{L^2d^2}{8}|\mu_{n+1} - \mu_n|^2 + \frac{L^2 d^2 \mu_1^2}{8}\right]  \\
        &\leq \frac{L^2}{8}  d^2\sum_{n=1}^\infty \frac{(\mu_n - \mu_{n+1})^2}{p_n} + \frac{L^2 }{8} d^2\mu_1^2.
    \end{align*}
    Therefore, we finally obtain
\begin{align*}
    \Var[\sfP_3(n,v) \, v ]  = & d\EE[(\sfP_3(n,v)- \sfP_4(n,v))^2  ] + \Var[\sfP_4(n,v)\, v] \\ 
    \leq & \Var[\sfP_4(n,v)\, v] + \frac{L^2}{8}  d^3\sum_{n=1}^\infty \frac{(\mu_n - \mu_{n+1})^2}{p_n} + \frac{L^2 }{8} d^3\mu_1^2.
\end{align*}
It concludes the proof. 
\end{proof}

% Recall that $\mathsf{P}_3(n,v)$ is defined as 
% \begin{equation}
% \mathsf{P}_3(n, v)= \frac{f(x+\mu_1 v)-f(x)}{\mu_1}\mathrm{U}_2 + \frac{1}{p_n}\Biggl[ \frac{f(x+\mu_{n+1} v)-f(x)}{\mu_{n+1}} - \frac{f(x+\mu_n v)-f(x)}{\mu_n} \Biggr](1-\mathrm{U}_2).  \nonumber
% \end{equation}
% where $\mathrm{U}_2\sim  \operatorname{Uniform}\left( \{ 0, 1 \}\right)$ is a selection variable.  Define a conjugate estimator
% \begin{equation}
% \widebar{\mathsf{P}}_3(n, v)= \frac{f(x+\mu_1 v)-f(x)}{\mu_1}(1-\mathrm{U}_2) + \frac{1}{p_n}\Biggl[ \frac{f(x+\mu_{n+1} v)-f(x)}{\mu_{n+1}} - \frac{f(x+\mu_n v)-f(x)}{\mu_n} \Biggr]\mathrm{U}_2.  \nonumber
% \end{equation}
% By symmetry, they have the same variance and the same expectation, and 
% $$\sfP_4(n,v) v =  {\mathsf{P}}_3(n, v)v + \widebar{\mathsf{P}}_3(n, v) v.$$
% We evaluate the variance on both sides:
% \begin{align*}
%     &\Var[\sfP_4(n,v) v] \\
%     =  &\Var[\widebar{\mathsf{P}}_3(n, v) v] + \Var[\sfP_3(n,v) v] + 2\EE[ \left( \widebar{\mathsf{P}}_3(n, v) v - \nabla f(x) \right)^\top \left(  {\mathsf{P}}_3(n, v) v - \nabla f(x) \right) ]. 
% \end{align*}
% As the result,
% \begin{align*}
%     \Var[\sfP_4(n,v) v] = 2 \Var[\sfP_3(n,v) v] - \|\nabla f(x) \|^2 +2d  \EE \left[\widebar{\mathsf{P}}_3(n, v)  {\mathsf{P}}_3(n, v)  \right].
% \end{align*}
% It suffices to evaluate $\EE \left[\widebar{\mathsf{P}}_3(n, v)  {\mathsf{P}}_3(n, v) \right]$. Denote
% $$\Delta_{n} = \delta_n(v) + \nabla f(x)^\top v = \frac{f(x+\mu_n v) - f(x)}{\mu_n}.$$
% Then we have
% \begin{align*}
%     \EE \left[\widebar{\mathsf{P}}_3(n, v)  {\mathsf{P}}_3(n, v) \right]  &=  \EE\left[ \mathrm{U}_2^2  \Delta_1 \frac{\Delta_{n+1} - \Delta_n}{p_n}\right]  + \EE\Biggl[ (1-\mathrm{U}_2)^2 \Delta_1 \frac{\Delta_{n+1} - \Delta_n}{p_n}\Biggr]  \\ 
%     &=  \EE\Biggl[ \Delta_1 \frac{\Delta_{n+1} - \Delta_n}{p_n}\Biggr] \\  
%     &\overset{(i)}{\leq} \frac{1}{2}\EE\Biggl[ \Delta_1  + \frac{\Delta_{n+1} - \Delta_n}{p_n}\Biggr]^2 \\
%     &\overset{(ii)}{\leq} \frac{1}{2}\EE  \sfP_4(n,v)^2  ,
% \end{align*}
% where (i) applies $ab\leq \frac{1}{2}(a+b)^2$ and (ii) applies the definition of $\sfP_4(n,v)$. As 
% $$ \Var[\mathsf P_4(n,v)\,v  ] = d \EE [\mathsf P_4(n,v)^2 ] - \|\nabla f(x) \|^2,$$
% we obtain 
% \begin{align*}
%     \Var[\sfP_4(n,v) v] &= 2 \Var[\sfP_3(n,v) v] - \|\nabla f(x) \|^2 +2d  \EE \left[\widebar{\mathsf{P}}_3(n, v)  {\mathsf{P}}_3(n, v)  \right] \\ 
%     &\geq   2 \Var[\sfP_3(n,v) v] - \|\nabla f(x) \|^2 - d \EE  \sfP_4(n,v)^2  \\ 
%     &=  2 \Var[\sfP_3(n,v) v] -2 \|\nabla f(x) \|^2 -\Var[\sfP_4(n,v) v].
% \end{align*}
% Therefore, we obtain 
% \begin{align*}
%    \Var[\sfP_3(n,v) v] \leq \Var[\sfP_4(n,v) v] + \|\nabla f(x) \|^2. 
% \end{align*}



% \newpage

% the variance of $\mathsf{P}_3(n, v) v$ is equal to $\EE[\mathsf{P}_3(n, v)^2 ] - \|\nabla f(x)\|^2$. Therefore, it suffices to consider $\EE[\mathsf{P}_3(n, v)^2 ]$. 

% Recall that $\mu_n = \frac{\mu}{n}$. 

% Variance 
% $$\frac{2}{\mu_1^2}+ \frac{2}{\mu_{n+1}^2 p_n^2} + \frac{2}{\mu_{n}^2 p_n^2}$$

% Because $\mathsf{P}_4(n, v)= \frac{f(x+\mu_1 v)-f(x)}{\mu_1} + \frac{1}{p_n}\Biggl[ \frac{f(x+\mu_{n+1} v)-f(x)}{\mu_{n+1}} - \frac{f(x+\mu_n v)-f(x)}{\mu_n} \Biggr]$, we consider the second term:
% \begin{align*}
%      &\frac{1}{p_n}\Biggl[ \frac{f(x+\mu_{n+1} v)-f(x)}{\mu_{n+1}} - \frac{f(x+\mu_n v)-f(x)}{\mu_n} \Biggr] \\ 
%      =&\frac{1}{p_n}\Biggl[ (n+1)\frac{f(x+\mu_{n+1} v)-f(x)}{\mu} - n \frac{f(x+\mu_n v)-f(x)}{\mu} \Biggr] \\ 
%      =& \frac{1}{\mu p_n}\Biggl[ n f(x+\mu_{n+1} v)- n  f(x+\mu_n v) + f(x+\mu_{n+1} v) -  f(x)     \Biggr]
% \end{align*}


% \begin{align*}
%      &\frac{1}{p_n}\Biggl[ \frac{f(x+\mu_{n+1} v)-f(x)}{\mu_{n+1}} - \frac{f(x+\mu_n v)-f(x)}{\mu_n} \Biggr] \\ 
%      \leq &\frac{1}{p_n}\Biggl[   \left( \mu_{n+1} - \mu_n  \right) \widebar{L}\Biggr]   
% \end{align*}
% COCO\todo{need a better upper bound here}



% \begin{align*}
%     \EE[\mathsf{P}_4(n, v)^2 ] &= \left( \frac{f(x+\mu v)-f(x)}{\mu} \right)^2  \\ 
%     &\quad +  \left(  \frac{1}{\mu p_n}\Biggl[ n f(x+\mu_{n+1} v)- n  f(x+\mu_n v)+ f(x+\mu_{n+1} v) -  f(x)     \Biggr]\right)^2 \\ 
%     &\quad + 2 \EE \left(  \frac{f(x+\mu v)-f(x)}{\mu}  \frac{1}{p_n}\Biggl[ \frac{f(x+\mu_{n+1} v)-f(x)}{\mu_{n+1}} - \frac{f(x+\mu_n v)-f(x)}{\mu_n} \Biggr] \right) 
% \end{align*}

% \begin{itemize}
%     \item Standard variance:
%     \begin{align*}
%         \left( \frac{f(x+\mu v)-f(x)}{\mu} \right)^2
%     \end{align*}

%     \item Known expectation:
%     \begin{align*}
%         &2 \EE \left(  \frac{f(x+\mu v)-f(x)}{\mu}  \frac{1}{p_n}\Biggl[ \frac{f(x+\mu_{n+1} v)-f(x)}{\mu_{n+1}} - \frac{f(x+\mu_n v)-f(x)}{\mu_n} \Biggr] \right) \\ 
%         =& 2 \EE \left(  \frac{f(x+\mu v)-f(x)}{\mu}  \left[  \nabla_v f(x) -  \frac{f(x+\mu v)-f(x)}{\mu}\right] \right) \\ 
%         =& 2 \EE \left(  \frac{f(x+\mu v)-f(x)}{\mu}    \nabla_v f(x)   \right) - 2 \EE \left(  \frac{f(x+\mu v)-f(x)}{\mu}    \right)^2 
%     \end{align*} \todo{not correct; evaluate the expectation again}

%     \item Key term: 
%     \begin{align*}
%          &\left(  \frac{1}{\mu p_n}\Biggl[ n f(x+\mu_{n+1} v)- n  f(x+\mu_n v)+ f(x+\mu_{n+1} v) -  f(x)     \Biggr]\right)^2 \\ 
%          \leq & \frac{1}{\mu^2 p_n^2} \left[ 2 \left( n f(x+\mu_{n+1} v)- n  f(x+\mu_n v) \right)^2 + 2 \left(f(x+\mu_{n+1} v) -  f(x)  \right)^2  \right] \\ 
%          \leq & \frac{1}{\mu^2 p_n^2} \left[ 2n^2 G^2  \left( \mu_{n+1} - \mu_n \right)^2 + 2 \left(f(x+\mu_{n+1} v) -  f(x)  \right)^2  \right]
%     \end{align*}

%     Rough estimate:
%     \begin{align*}
%         \frac{(\mu_{n+1} - \mu_n)^2}{p_n^2}C
%     \end{align*}
% \end{itemize}
% Put together:
% \begin{align*}
%     \Var &\leq  2 \EE \left(  \frac{f(x+\mu v)-f(x)}{\mu}    \nabla_v f(x)   \right) -   \EE \left(  \frac{f(x+\mu v)-f(x)}{\mu}  \right)^2 +  \frac{(\mu_{n+1} - \mu_n)^2}{p_n^2}C \\
%     &= \EE \left[ \frac{f(x+\mu v)-f(x)}{\mu} 
%  - \nabla_v f(x)\right]^2  -\EE\|\nabla_v f(x)\|^2 +  \frac{(\mu_{n+1} - \mu_n)^2}{p_n^2}C \\ 
%  &= \mu^2 C^2 - \EE \| \nabla f(x) \|^2  +  \sum_{n=1}^\infty \frac{(\mu_{n+1} - \mu_n)^2}{p_n }C 
% \end{align*}


% $\mu_n = \alpha^n \mu$. Then $(\mu_{n+1} - \mu_n)^2 = \alpha^{2n} \mu^2 (\alpha - 1)^2$ 

% If consider variance:
% \begin{align*}
%      \sum_n \frac{C}{\mu_{n+1}^2 + \mu_n^2 } \frac{1}{p_n} + \frac{(\mu_{n+1} - \mu_n)^2}{p_n }D
% \end{align*}

  


\subsection{On the Optimal Sampling Distribution and the Perturbation Stepsize Sequence}\label{appendix:optimal-sampling}
We recap the full statement of \Cref{prop:lower_bound_varrho}. 

\begin{theorem}\label{prop:lower_bound_varrho-appendix}
Let $\{\mu_n\}_{n=1}^\infty$ be a positive, decreasing sequence with   $\sum_{n=1}^\infty\mu_n<\infty$, and let $\{p_n\}_{n=1}^\infty$ be a PMF.   Then the following statements hold:
\begin{enumerate}[label=(\alph*)]
    \item Define $\varrho_n \;=\; \frac{(\mu_{n+1}-\mu_n)^2}{p_n}.$   The lower bound of $\varrho$ is given by 
    $$\varrho = \sum_{n=1}^\infty \varrho_n \;\ge\; \mu_1^2.$$  
    Moreover, equality holds if and only if
$p_n = \frac{\mu_n - \mu_{n+1}}{\mu_1}.$ 
    \item  Define $\varphi_n \;=\; \frac{ \mu_n^2}{p_n}.$   The lower bound of $\varphi$ is given by 
    $$\varphi= \sum_{n=1}^\infty \varphi_n \geq \Big( \sum_{n=1}^\infty \mu_n \Big)^2> \mu_1^2.$$
    Moreover, equality holds if and only if
$p_n = \frac{\mu_n }{\sum_{n=1}^\infty \mu_n }.$ 
\end{enumerate}
\end{theorem} 
\begin{proof}
Write $a_n=\mu_{n+1}-\mu_n$, so $\sum_{n=1}^\infty a_n=\mu_1$  and $\sum_n p_n=1$.  By Cauchy-Schwarz inequality,
\[
\Bigl(\sum_n |a_n|\Bigr)^2
\le
\Bigl(\sum_n p_n\Bigr)
\Bigl(\sum_n \frac{a_n^2}{p_n}\Bigr)
=
\sum_{n=1}^\infty \varrho_n,
\]
which yields the claimed lower bound in (a).  Equality occurs exactly when $p_n\propto|a_n|$, i.e.\ $p_n=(\mu_n-\mu_{n+1})/\mu_1$. Similarly,  By Cauchy-Schwarz inequality, 
$$\Bigl( \sum_n \mu_n  \Bigr)^2 \leq\Bigl( \sum_n  \frac{\mu^2_n}{p_n} \Bigr) \Bigl( \sum_n p_n  \Bigl) = \sum_{n=1}^\infty \varphi_n,$$
which yields the claimed lower bound in (b). 
\end{proof} 

In the following example, we include the omitted details of \Cref{example:zipf}. 
\begin{example} \label{example:zipf-appendix}

We consider the Zipf distribution $n \sim \operatorname{Zipf}(s)$ ($s>1$).  Then
\[
p_n=\frac{1}{\zeta(s)}\frac{1}{n^s}, 
\qquad
\zeta(s)=\sum_{n=1}^\infty \frac{1}{n^s}.
\]
We define $\{\mu_n\}$ by the recursion
\[
\mu_n - \mu_{n+1} \;=\;\mu_1\,p_n
\;=\;\mu_1\,\frac{1}{\zeta(s)}\frac{1}{n^s},
\]
so that summing gives the closed‐form
\[
\mu_n
=\mu_1\Bigl(1-\frac{\sum_{j=1}^{n-1}j^{-s}}{\zeta(s)}\Bigr)
=\mu_1\,\frac{\sum_{j=n}^\infty j^{-s}}{\zeta(s)}.
\]
A direct check shows this choice attains the lower bound $\varrho=\mu_1^2$ on $\sum(\mu_n-\mu_{n+1})^2/p_n$. Now we turn to bound $\varphi$: 
\[
\varphi
=\sum_{n=1}^\infty \frac{\mu_n^2}{p_n}
=\mu_1^2\,\zeta(s)\sum_{n=1}^\infty n^s
\Bigl(1-\tfrac{\sum_{j=1}^{n-1}j^{-s}}{\zeta(s)}\Bigr)^2
=\mu_1^2\,\zeta(s)\sum_{n=1}^\infty n^s
\Bigl(\tfrac{\sum_{j=n}^\infty j^{-s}}{\zeta(s)}\Bigr)^2.
\]
For $s>3$, use the integral bound
\[
\sum_{j=n}^\infty j^{-s}
\;\le\;\int_{n-1}^\infty x^{-s}\,dx
=\frac{(n-1)^{1-s}}{s-1},
\]
to get
\[
\frac{\mu_n^2}{p_n}
\;\le\;\frac{\mu_1^2}{(s-1)^2\,\zeta(s)}\;n^{2-s}.
\]
Since $s>3$ the series $\sum_{n=1}^\infty n^{2-s}=\zeta(s-2)$ converges, giving the clean bound
\[\varphi
\;=\;\sum_{n=1}^\infty\frac{\mu_n^2}{p_n}
\;\le\;\frac{\zeta(s-2)}{(s-1)^2\,\zeta(s)}\;\mu_1^2. 
\] 
\end{example}


\subsection{Discussions: Variance of Random Directional Derivative} 
In this subsection, we analyze the variance of a gradient estimate based on a random directional derivative. Let \( v \) be a random vector uniformly sampled from the sphere with the dimension $d$. We approximate the gradient \( \nabla f(x) \) using the random directional derivative defined as  
\[
\nabla_v f(x) := vv^\top \nabla f(x).
\]  
Assuming that the expectation satisfies \( \mathbb{E}[vv^\top] = I \), it is essential to evaluate the variance of this estimator. Specifically, we compute:  
\begin{align*}
    \mathbb{E} [\nabla f(x)^\top vv^\top   vv^\top \nabla f(x) ] 
    &= \mathbb{E} [\|\nabla f(x) \|^2 \|v\|^2] \\ 
    &= d \|\nabla f(x) \|^2. 
\end{align*}  
Thus, the variance is given by  
\[
\text{Var}[\nabla_v f(x)] = d \|\nabla f(x) \|^2 - \|\nabla f(x) \|^2 = (d-1)\|\nabla f(x) \|^2.
\]  

This result indicates that even when the exact directional derivative is available, the variance still scales with \(\mathcal{O}(d)\) relative to the gradient norm. Consequently, it is not avoidable to remove the dependence on the dimension $d$.



\section{Convergence Analysis}
In this section, we present the proof of \Cref{cor:sgd-convergence}.  Recall that our goal is to solve the stochastic optimization problem \Cref{eq:optimization-problem}:
\begin{align*} 
\min_{x \in \RR^d} f(x) := \EE_{\xi \sim \Xi} f(x;\xi),
\end{align*}
where the second-order continuously differentiable function $f(\cdot;\xi):\RR^d \to \RR$ has $L$-Lipschitz gradient for every $\xi$. We consider the convergence upper bound of the classical stochastic gradient descent (SGD) algorithm with the constant learning rate $\eta$ given the initialization $x_0$:
\begin{align}\label{eq:sgd-iteration}
   x_{t+1} = x_t - \eta g(x_t), \tag{SGD}
\end{align}
where $g(x_t)$ is an unbiased estimator of the full gradient $\nabla f(x)$. 

\subsection{The Convergence Upper Bound}\label{sec:convergence-proof-appendix}
The full statement of \Cref{cor:sgd-convergence} is given as follow:
\begin{corollary}
Consider the stochastic optimization problem in \Cref{eq:optimization-problem}, and suppose that the individual loss \( f(x;\xi) \) is second-order differentiable with \( L \)-Lipschitz continuous gradient in \( x \), uniformly over \(\xi\sim\Xi\). Assume the stochastic gradient is approximated using the $\sfP_k$-estimator \( \sfP_k(n,v)\,v \) for \(k =2,3,4\). Let the SGD iteration be defined as 
$$x_{t+1} = x_t - \eta \sfP_k(n_t, v_t)\,v_t$$ where $\eta \in (0, \frac{1}{L^2d}]$ is the stepsize. Then the iterates satisfy the following convergence guarantee:
    $$\min_{0\leq t \leq T-1} \EE \|\nabla f(x_t)\|^2 \leq L  \eta \left[ C_k d^3 \mu^2 + 2dL (f^* - \EE_{\xi\sim\Xi} f^*_\xi)\right] + \frac{2}{\eta T}\underline{\delta},$$
    where $\underline{\delta}:=f(x_0) - f^*$ and $C_k =\begin{cases}
        \quad \frac{28L^2}{12} &k=2, \\ 
        \quad \frac{3L^2}{4}   &k=3, \\
        \quad \frac{5L^2}{4}  &k=4, \\
    \end{cases}$ ($k=2,3,4$) is the estimator-dependent error term. Consequently, choosing \(\eta = \Theta(1/\sqrt{dT})\) and $\mu = \mathcal{O}(\frac{1}{d})$ yields the optimal complexity $T = \Theta(\frac{d}{\epsilon^4})$ of having $\min_{0\leq t \leq T-1} \EE \|\nabla f(x_t)\| \leq \epsilon$. 
\end{corollary}
\begin{proof}
    Let $g(x) = \sfP_k(n,v)\,v$ for $k=2,3,4$. By \Cref{lemma:2}, we have
    $$\EE \|g(x) \| \leq d L \|\nabla f(x)\|^2 +  C_k d^3 \mu^2 + 2dL (f^* - \EE_{\xi\sim\Xi} f^*_\xi).$$
    Then we set $B=d L$ and $C= C_k d^3 \mu^2 + 2dL (f^* - \EE_{\xi\sim\Xi} f^*_\xi)$ in \Cref{lemma:1}.  As the result, when $\eta \leq \frac{1}{L^2 d}$,  we have 
    $$\min_{0\leq t \leq T-1} \EE \|\nabla f(x_t)\|^2 \leq L  \eta \left[ C_k d^3 \mu^2 + 2dL (f^* - \EE_{\xi\sim\Xi} f^*_\xi)\right] + \frac{2}{\eta T}\underline{\delta}.$$
    The complexity of making $\min_{0\leq t \leq T-1} \EE \|\nabla f(x_t)\| \leq \epsilon$ is given by
    $$T \geq \frac{12 \underline{\delta} L}{\epsilon^2} \max \left\{ d L, 2\frac{ C_k d^3 \mu^2 + 2dL (f^* - \EE_{\xi\sim\Xi} f^*_\xi) }{ \epsilon^2} \right\}.$$
    Setting $\mu = \calO(\frac{1}{d})$, the optimal complexity is given by $T=\Theta(\frac{d}{\epsilon^4})$.
\end{proof}


\subsection{Supporting Lemmas}
\begin{lemma}\label{lemma:0} 
    Let $f^*:=\min_{x} f(x)$ and $f^*_\xi:=\min_{x} f(x;\xi)$. If for each $\xi$,  $f(x;\xi)$ has $L$-Lipschitz gradient, then
    $$\EE \| \nabla f(x;\xi)\|^2 \leq L \|\nabla f(x)\|^2 + 2L (f^* - \EE_{\xi\sim\Xi} f^*_\xi).$$
\end{lemma}
\begin{proof}
    This lemma is directly taken from Proposition 2, \citet{khaled2022better}. 
\end{proof}


\begin{lemma}\label{lemma:1}
    Let the second-order continuously differentiable function $f(\cdot\; ;\xi):\RR^d\to \RR$ be lower bounded by $f^*:=\min_x f(x)$ and have $L$-Lipschitz gradient  for every $\xi$, and $g(x)$ be an unbiased estimator of $\nabla f(x)$. Suppose that $g(x)$ and $f(x)$  satisfy the expected smoothness condition
    $$\EE \|g (x) \|^2 \leq   B\cdot \|\nabla f(x)\|^2 + C.$$
    If the learning rate $\eta\leq \frac{1}{LB}$ and define $\underline{\delta}:=f(x_0)-f^*$, then the $T$-th iteration of \ref{eq:sgd-iteration} satisfies
    $$\min_{0\leq t \leq T-1} \EE \|\nabla f(x_t)\|^2 \leq LC \eta + \frac{2}{\eta T}\underline{\delta}.$$
\end{lemma}
\begin{proof}
    This lemma is directly taken from Theorem 2, \citet{khaled2022better}.
\end{proof}

\begin{lemma}\label{lemma:2}
    Let $g(x)$ be the $\sfP_k$-estimator (for $k=2,3,4$), $f^*:=\min_{x} f(x)$, and $f^*_\xi:=\min_{x} f(x;\xi)$. Then the expected smoothness condition is given by
    $$\EE \|g(x) \| \leq d L \|\nabla f(x)\|^2 +  C_k d^3 \mu^2 + 2dL (f^* - \EE_{\xi\sim\Xi} f^*_\xi).$$
    where $$C_k =\begin{cases}
        \quad \frac{28L^2}{12} &k=2, \\ 
        \quad \frac{3L^2}{4}   &k=3, \\
        \quad \frac{5L^2}{4}  &k=4, \\
    \end{cases}$$
    is the error term introduced by the zeroth-order gradient estimation. For $\sfP_2$-estimator, we choose $\{\mu_n\}$ and $\{p_n\}$ such that $\varrho=\mu^2$ and $\varphi\leq 2\mu^2$. 
\end{lemma}
\begin{proof}
    Let $\hat{\nabla} f(x;\xi)$ be the $\sfP_k$-estimator for $k=2,3,4$. First, we notice the following variance-decomposition holds: 
    \begin{align*}
        \EE \| \hat{\nabla} f(x;\xi) \|^2 &=  \EE \| \hat{\nabla} f(x;\xi)  - \nabla f(x;\xi) + \nabla f(x;\xi)\|^2 \\ 
        &= \EE \| \nabla f(x;\xi) \|^2 + \EE \left[ \Var\left[  \sfP_k \ v \mid \xi\right] \right].
    \end{align*}
    By \Cref{thm:variance} and \Cref{prop:lower_bound_varrho}, we set 
    $$\Var\left[  \sfP_k \ v \mid \xi\right]  \leq (d-1) \|\nabla f(x;\xi) \|^2 + C_k d^3 \mu^2,$$
    where $C_k$ is determined by the estimator; we also assume that the optimal distribution and perturbations are taken obeying \Cref{prop:lower_bound_varrho}. As the result,
    \begin{align*}
        \EE \| \hat{\nabla} f(x;\xi) \|^2  &\leq  \EE \| \nabla f(x;\xi) \|^2 +  + \EE \left[ (d-1) \|\nabla f(x;\xi) \|^2 + C_k d^3 \mu^2\right] \\ 
        &\leq d \EE \| \nabla f(x;\xi) \|^2 +  C_k d^3 \mu^2 \\ 
        &\overset{(i)}{\leq }d L \|\nabla f(x)\|^2 +  C_k d^3 \mu^2 + 2dL (f^* - \EE_{\xi\sim\Xi} f^*_\xi),
    \end{align*}
    where (i) applies \Cref{lemma:0}. It completes the proof.
\end{proof}

% \begin{lemma}
%     Let the second-order continuously differentiable function $f:\RR^d\to \RR$ be lower bounded by $f_*$ and have $L$-Lipschitz gradient, and $g(x)$ be an unbiased estimator of $\nabla f(x)$. Suppose that (1) $g(x)$ and $f(x)$  satisfy the expected smoothness condition
%     $$\EE \|g (x) \|^2 \leq B\cdot \|\nabla f(x)\|^2 + C;$$
%     and (2) the function $f:\RR^d\to \RR$ satisfies the  Polya-Łojasiewicz condition 
%     $$\frac{1}{2}\|\nabla f(x)\|^2 \geq \mathscr{c} \left( f(x) - f_*\right)$$
%     for all $x\in \RR^d$. If the learning rate $\eta\leq \frac{1}{2LB}$, then the $T$-th iteration of \ref{eq:sgd-iteration} satisfies
%     $$ \EE \left[ f(x_T) - f_*\right] \leq $$
% \end{lemma}
% \begin{proof}
%     Theorem 3 from \citet{khaled2022better}.
% \end{proof}




\section{Experiments Details}\label{appendix:experiments}
In this section, we describe the detailed experiment setting and the hyperparameter configurations.
\subsection{Synthetic Example} 
In the synthetic example, we compared gradient estimators across varying dimensions ($d = 16, 64, 256, 1024, 4096$) using both quadratic and logistic functions. For fair comparison, all estimators used a consistent number of function evaluations of $3$ and the perturbation stepsize $\mu = 10^{-5}$ (for $\sfP_3$-estimator, we set $\mu_1= \mu$). The $\sfP_3$-estimator was configured with parameter $s = 2.0$ and followed the same perturbation stepsize scheduling as \Cref{example:zipf}. We evaluated four gradient estimators: Zipf's $P_3$-estimator, one-side two-point estimator with the Gaussian smoothing, one-side two-point estimator with the uniform smoothing, and two-side two-point estimator with the Gaussian smoothing. Each configuration was tested over $100$ independent trials with a fixed random seed for reproducibility.  

\paragraph{Code Availability and System Requirements} 
All source code is included in the supplementary materials. No specific hardware is required; any machine supporting Python 3.10.10 should suffice. A Jupyter notebook version is also provided for convenient execution on Google Colab.

\subsection{Language Model Optimization}
In the language model optimization experiment, we compare the performance of different gradient estimators within a vanilla SGD framework for fine-tuning a language model on a sentiment classification task. The learning rate of SGD is fixed at $\eta = 10^{-4}$, the batch size of SGD is fixed at $16$ (this batch size corresponds to the number of stochastic samples and is different from the batch size in multiple-point zeroth-order estimator), and the perturbation stepsize is set to $\mu = 10^{-3}$ (for our proposed unbiased estimators, we use $\mu_1 = \mu$). Due to limitations in numerical precision, we do not sample the extreme tail of the distribution. Instead, we truncate the sampling distribution by enforcing $p_n \geq 10^{-3}$ to ensure numerical stability and avoid excessively small probabilities. All remaining hyperparameters are summarized in \Cref{tab:hyper}.


\begin{table}[t]
\centering
\caption{Summary of gradient estimators used in the language model optimization experiment.}\label{tab:hyper}
\renewcommand{\arraystretch}{1.2}  % Slightly increase row height
\resizebox{\textwidth}{!}{%
\begin{tabular}{lcccc}
\toprule
\textbf{Estimation Method} & \textbf{Estimator Formula} & \textbf{Batch Size} $b$ & \textbf{\# Function Calls} & \textbf{Notes} \\
\midrule
One-Side Two-Point Estimator & $\displaystyle\sum_{i=1}^b\dfrac{f(x + \mu v_i) - f(x)}{\mu} v_i$ & 3 & 4 & $v_i \SimIID  \operatorname{Normal}(0, I_d)$ \\
Two-Side Two-Point Estimator & $\displaystyle\sum_{i=1}^b\dfrac{f(x + \mu v_i) - f(x-\mu v_i)}{2\mu} v_i$ & 2 & 4 & $v_i \SimIID \operatorname{Normal}(0, I_d)$ \\ 
\midrule
Zipf's $\sfP_3$-Estimator & \Cref{eq:P3_def} & 1 & 3 & $s=1.5$, \Cref{example:zipf} \\
Geometric $\sfP_3$-Estimator & \Cref{eq:P3_def} & 1 & 3 & $c=0.5$, \Cref{example:geometric} \\
Zipf's $\sfP_4$-Estimator & \Cref{eq:P4_def} & 1 & 4 & $s=1.5$, \Cref{example:zipf} \\
Geometric $\sfP_4$-Estimator & \Cref{eq:P4_def} & 1 & 4 & $c=0.5$, \Cref{example:geometric} \\
\bottomrule
\end{tabular}
}
\label{tab:gradient-estimators}
\end{table}
 



\paragraph{Code Availability and System Requirements}
All source code and reproduction instructions are provided in the supplementary materials. Experiments were conducted on a cluster running RHEL 8, equipped with dual AMD EPYC 7763 processors, 512 GB of memory, and seven NVIDIA RTX 5000 GPUs. To reproduce the language model optimization experiments, we recommend using a CUDA-compatible GPU with at least 24 GB of VRAM.



\section{Broader Impact}\label{appendix:braoder-impact}
This work introduces a new class of unbiased zeroth-order gradient estimators that offer both theoretical guarantees and practical advantages. By eliminating bias without increasing variance, our method enhances the reliability of optimization in settings where gradient information is unavailable or costly to obtain. These include fine-tuning large language models under memory constraints, conducting black-box adversarial robustness evaluations, and solving scientific computing tasks such as physics-informed neural networks. On the theoretical side, our estimators advance the understanding of zeroth-order optimization and provide new tools for the zeroth-order gradient estimation. This work opens a promising direction for future research in gradient-free optimization and its broad applications in machine learning and beyond.
 
\section{Limitations}\label{appendix:limitations}
Despite the theoretical guarantees and empirical improvements demonstrated by our proposed unbiased zeroth-order gradient estimators, several limitations remain. First, the estimators rely on sampling from an infinite sequence of perturbation steps, but practical implementations must truncate this sequence, which may reintroduce bias or affect variance control. Second, the proposed estimators are based on directional derivatives, which inherently exhibit a dependence on the problem dimension $d$; this dimensional dependence is generally unavoidable for this class of methods. Lastly, while we validate the approach on synthetic tasks and language model fine-tuning, we have not extensively evaluated its performance across a broader range of optimization problems, and the observed empirical gains may not fully generalize to settings involving non-smooth objectives or high levels of evaluation noise.

%%%%%%%%%%%%%%%%%%%%%%%%%%%%%%%%%%%%%%%%%%%%%%%%%%%%%%%%%%%%

\end{document}
